\documentclass[a4paper,12pt]{ctexart}
\usepackage{xcolor}
\usepackage{graphicx}
\usepackage[top=1.8cm, bottom=1.8cm, left=1.8cm, right=1.8cm]{geometry}
\usepackage{array}
\usepackage{hyperref}
\hypersetup{colorlinks=true,linkcolor=blue!70!black}
\linespread{1.35}

\definecolor{storyblue}{RGB}{0,102,153}
\definecolor{thinkorange}{RGB}{230,120,20}
\definecolor{funboxbg}{RGB}{245,248,252}

\newenvironment{thinkbox}{
  \vspace{0.3cm}
  \begin{center}
  \begin{minipage}{0.88\textwidth}
  \colorbox{funboxbg}{
  \begin{minipage}{0.94\textwidth}
  \vspace{0.15cm}
  \textbf{\textcolor{thinkorange}{【 想一想 】}}
}{
  \vspace{0.15cm}
  \end{minipage}}
  \end{minipage}
  \end{center}
  \vspace{0.3cm}
}

\newenvironment{funfact}{
  \vspace{0.2cm}
  \begin{center}
  \begin{minipage}{0.88\textwidth}
  \fbox{
  \begin{minipage}{0.92\textwidth}
  \vspace{0.1cm}
  \textbf{\textcolor{storyblue}{【 有趣的小知识 】}}
}{
  \vspace{0.1cm}
  \end{minipage}}
  \end{minipage}
  \end{center}
  \vspace{0.2cm}
}

\newenvironment{figplaceholder}[1]{
  \vspace{0.2cm}
  \begin{center}
  \fbox{
  \begin{minipage}{0.82\textwidth}
  \centering
  \textcolor{blue!70!black}{\textbf{【插图位置】}} \\[0.2cm]
  #1
  \end{minipage}}
  \end{center}
  \vspace{0.2cm}
}

\title{\textcolor{blue!70!black}{\Huge\textbf{自然科学小故事}}}
\author{}
\date{}

\begin{document}
\maketitle
\thispagestyle{empty}

\section*{\textcolor{blue!70!black}{给小朋友的话}}

这个世界充满了神奇的现象——为什么天是蓝的？为什么苹果会掉下来？为什么小船能浮在水面上？这个模块不讲公式、不背概念，只讲故事。每一个故事背后，都藏着一个有趣的自然科学道理。

你不需要一下子把所有道理都弄明白。你只需要带着好奇心，像一位小侦探一样，跟着故事里那些了不起的科学家们，一起去发现世界的秘密。

准备好了吗？让我们一起开始这场奇妙的自然探索之旅吧！

\vspace{0.5cm}

\newpage

\section{\textcolor{orange!80!black}{第一课\quad 牛顿与苹果——万有引力}}

\subsection*{树下的年轻人}

1666年，距离现在大约360年。在英国一个叫伍尔索普的小村庄里，住着一个23岁的年轻人，他的名字叫\textbf{艾萨克·牛顿}。

那一年，伦敦爆发了一场可怕的瘟疫，大学都停课了。牛顿只好回到乡下的老家，在自家果园里看书、思考问题。他每天都在想一些看起来毫无用处的问题，比如：月亮为什么会绕着地球转？东西为什么总是往下掉，而不是往上飞？

有一天，牛顿正坐在一棵苹果树下想得出神。突然——"咚"的一声，一颗苹果从树上掉了下来，正好落在他脚边。

换作是我们，大概会说："哦，苹果熟了。"然后——也许——捡起来咬一口。

但牛顿没有。他盯着那颗苹果，脑子里冒出了一个了不起的问题。

\begin{thinkbox}
如果苹果会往下掉，那月亮为什么不会掉下来呢？
\end{thinkbox}

\subsection*{一个改变世界的"为什么"}

这个问题听起来有点傻，对不对？苹果是苹果，月亮是月亮，它们怎么能一样呢？

但牛顿不这么想。他反复琢磨：苹果从树上掉下来，是因为地球在"拉"它。那么，高高挂在天上的月亮，是不是也被地球"拉"着？如果是这样，月亮为什么没有砸到我们头上呢？

牛顿想了很久，终于想通了。他发现：地球确实在"拉"着月亮！但月亮之所以不会砸下来，是因为它一直在高速地往前"跑"。就像你拿一根绳子拴着一个小球转圈甩——只要速度够快，小球就会一直绕着你的手转，而不会飞走，也不会被绳子拉回来。

这里的"绳子"，就是地球和月亮之间肉眼看不见的吸引力。牛顿给它起了一个名字——\textbf{万有引力}。

\begin{figplaceholder}
此处插入牛顿坐在苹果树下的场景插图（见 figure/requirements/nature-intro-ch1.txt 插图1）
\end{figplaceholder}

\subsection*{万有引力到底是什么？}

简单来说，万有引力就是：\textbf{任何两个有质量的东西，都会互相吸引。}你和我之间、你和桌子之间、地球和月亮之间——都存在万有引力。

那你可能要问了：既然我和桌子互相吸引，为什么我没有被桌子吸过去呢？

问得好！这是因为牛顿还发现了一个重要的规律：引力的大小，和物体的\textbf{质量}有关。地球的质量太大了，所以它的引力能把你牢牢地"抓"在地面上；而桌子的质量太小了，它对你的引力微弱到你根本感觉不到。

正是因为地球有这么大的引力，我们才能稳稳地站在地上，而不是像气球一样飘到天上去。也是因为太阳的引力，地球才能年复一年地绕着太阳转，从不会"脱轨"跑到宇宙深处去。

\begin{funfact}
牛顿发现万有引力的时候，才23岁！但他怕自己的理论不够完善，足足等了20年才正式发表。所以啊，伟大的发现不仅需要聪明才智，还需要认真和耐心。
\end{funfact}

\begin{thinkbox}
跳起来的时候，你其实是在和整个地球"拔河"——你跳起来了，但地球的引力马上把你拉回来了。想想看：如果地球没有引力，跳一下会发生什么？
\end{thinkbox}

\begin{figplaceholder}
此处插入万有引力示意图：地球"拉"着月亮的简图（见 figure/requirements/nature-intro-ch1.txt 插图2）
\end{figplaceholder}

\subsection*{从苹果到火箭——引力与航天}

牛顿发现了万有引力，但他可能没有想到，几百年后，人类竟然利用这个原理，\textbf{反过来"对抗"了地球的引力}，飞出了地球！

这是怎么做到的呢？关键就在于牛顿发现的那个道理：\textbf{要想不被地球拉回来，你得跑得足够快}。

还记得我们说的"绳子拴小球"吗？如果你把球甩得足够快，球就能一直绕着你的手转。同样的道理，如果一艘飞船飞行的速度达到每秒7.9公里——这个速度叫做\textbf{第一宇宙速度}——它就能像月亮一样，绕着地球转圈而不会掉下来。我们熟悉的\textbf{天宫空间站}、\textbf{北斗卫星}，都是用这个速度在地球轨道上运行的。

如果想飞得更远——比如去月球、去火星——那就得再快一些。当速度达到每秒11.2公里（\textbf{第二宇宙速度}），飞船就能彻底"挣脱"地球的引力，飞向更遥远的太空。

你可能会好奇：那火箭是怎么飞起来的呢？

火箭的原理，其实很像一个被吹满了气然后松手的气球——气球里的气往后冲，气球就往前飞。火箭的发动机向下喷出炽热的气体，火箭就被推向反方向——也就是向上，冲向天空！这个原理叫\textbf{反冲运动}，也是牛顿发现的（牛顿第三定律：作用力和反作用力）。

所以你看，从一颗苹果掉在地上，到人类把探测器送到火星表面，背后是同一个道理在起作用。牛顿在苹果树下迈出的那一小步，最终带人类跨出了地球的一大步。

\begin{thinkbox}
中国的"嫦娥"探测器是怎么飞到月球的？它先被火箭送到地球轨道上（第一宇宙速度），然后再加速到第二宇宙速度，脱离地球引力飞向月球。到了月球附近，月球的引力又会把它"拉"住，让它绕着月球转——就像地球拉着月亮一样。是不是很巧妙？
\end{thinkbox}

\begin{funfact}
宇航员在空间站里会"飘"起来，不是因为那里没有引力！实际上，在空间站的轨道高度（大约400公里），地球的引力仍然有地面的90\%左右。宇航员之所以飘，是因为空间站一直在高速地"往下掉"——但它掉落的弧度和地球弯曲的表面正好"匹配"，所以永远落不到地上。这就是\textbf{失重}的真正原因，它和"没有引力"完全是两回事！
\end{funfact}

\begin{figplaceholder}
此处插入火箭发射与地球轨道示意图：展示从地面发射到进入轨道的过程（见 figure/requirements/nature-intro-ch1.txt 插图7）
\end{figplaceholder}

\newpage

\section{\textcolor{orange!80!black}{第二课\quad 阿基米德的浴缸——浮力的秘密}}

\subsection*{国王的难题}

两千多年前，在地中海的一个叫\textbf{叙拉古}的城邦里，住着古希腊最聪明的人之一——\textbf{阿基米德}。

有一回，叙拉古的国王打造了一顶纯金的新王冠。可是，国王总觉得金匠在金子里面掺了便宜的银子，偷了他的黄金。但他既没有证据，又舍不得把王冠切开检查——万一是自己冤枉了金匠呢？

这可是一道超级难题！怎么才能不破坏王冠，又检查出它是不是纯金的呢？国王把这个难题交给了阿基米德。

阿基米德想啊想，吃饭想，走路想，连洗澡都在想。

\subsection*{浴缸里的"尤里卡！"}

这一天，阿基米德拖着疲惫的身体进了浴室。他跨进盛满热水的浴缸——"哗啦"一声，水从浴缸边缘漫了出来，流了一地。

就在这一瞬间，阿基米德脑子里像划过一道闪电！他发现：自己的身体沉进水里，有一部分水被"挤"了出来。而且，身体沉下去越多，溢出来的水就越多！

阿基米德兴奋极了，他光着身子就从浴缸里跳了出来，一边跑上大街，一边大喊："\textbf{尤里卡！尤里卡！}"——这在希腊语里是"我发现了！我发现了！"的意思。

\begin{figplaceholder}
此处插入阿基米德在浴缸中发现浮力原理的场景插画（见 figure/requirements/nature-intro-ch1.txt 插图3）
\end{figplaceholder}

\subsection*{浮力的秘密}

阿基米德发现了什么呢？他发现了一个可以流传几千年的原理，后来人们把它称为\textbf{阿基米德原理}：

\textbf{任何物体放进水里，都会受到一个向上的托力，这个托力的大小，正好等于物体排开的水的重量。}

这个"向上的托力"，就是我们说的\textbf{浮力}。

回到国王的王冠。阿基米德想：金子和银子的密度不一样——同样重量的金子和银子，金子占的体积更小。如果把王冠浸入水中，它如果是纯金的，应该排开一定量的水；如果掺了银子，体积会更大，排开的水就会更多。

于是，阿基米德把王冠和一块同样重量的纯金分别放入水中，比较它们排出的水量。结果发现——王冠排出的水确实比纯金多！金匠果然掺假了。

国王的难题，就这样被一个浴缸解决了。

\begin{thinkbox}
拿两团一样重的橡皮泥，一团捏成球，一团捏成小船的形状，放进水里。谁会浮起来？为什么一样重的东西，一个沉一个浮呢？\\
\textcolor{gray}{\small 提示：这和排开水的多少有关——船的形状能排开更多的水，所以受到更大的浮力！}
\end{thinkbox}

\begin{funfact}
你有没有发现，在游泳的时候，你在水里感觉自己变"轻"了？这就是浮力在帮忙！水把你的身体往上托，所以你在水里能轻松地漂起来。大海里的浮力比游泳池更大哦，因为海水里溶解了盐，密度更大。这就是为什么在死海（一个特别特别咸的湖）里，人甚至可以躺在水面上看书而不会沉下去！
\end{funfact}

\begin{figplaceholder}
此处插入浮力原理示意图：物体入水排开水量的对比简图（见 figure/requirements/nature-intro-ch1.txt 插图4）
\end{figplaceholder}

\subsection*{轮船为什么不会沉？}

阿基米德原理不仅适用于浴缸和王冠，它还是所有船舶设计的\textbf{核心秘密}。

你有没有想过一个问题：一艘几十万吨的超级油轮，全是用钢铁造的，钢铁明明比水重得多，为什么它能稳稳地浮在海面上？

答案就在"排开水的体积"里。

造船的工程师把船底设计成巨大的、中空的形状——像一个巨大的铁碗扣在水面上。当船体沉入水中，它排开的水量非常非常大，因此受到的水的浮力也非常非常大。只要浮力等于整艘船的重力，船就能浮在水面上。

那为什么有时候船会沉呢？比如泰坦尼克号撞上冰山后，海水涌进了船体内部。船变重了，但排开的水量没有增加——浮力不够了，船就沉下去了。所以船底一定要保持密封，这是性命攸关的事。

\begin{figplaceholder}
此处插入轮船浮力原理示意图：铁块沉水 vs 铁船浮水对比（见 figure/requirements/nature-intro-ch1.txt 插图8）
\end{figplaceholder}

\subsection*{飞机又是怎么"浮"在空中的？}

你可能注意到了一个问题：飞机既不在水里，也没有排开水——它是怎么"浮"在空中的呢？

好问题！飞机的"浮力"和水的浮力完全不同，但我们可以把它理解成\textbf{空气给飞机的"托力"}。专业的叫法是\textbf{升力}。

如果你仔细观察飞机的机翼（翅膀），会发现它的形状很特别：上面是鼓起来的弧形，下面比较平。当飞机高速向前飞行时，空气从机翼上方和下方同时流过。因为上方是弧形的，空气要绕更远的路，所以流得更快；下方比较平，空气流得相对慢。

这里有一个神奇的规律（它叫\textbf{伯努利原理}）：\textbf{流体流动越快的地方，压力越小。}所以机翼上方的空气压力小，下方的压力大——下面的空气就会把机翼往上"托"。飞机就是靠这个"托力"飞起来的。

所以你看，虽然原理不同，但思路是一样的：\textbf{轮船靠水的浮力托起来，飞机靠空气的升力托起来。}一个向下排水，一个向上"托"气——异曲同工！

\begin{thinkbox}
放风筝的时候，风从风筝下面吹过，把风筝托上天空。这和飞机翅膀产生升力的道理有点像哦！下次放风筝的时候，可以感受一下那股向上的力量。
\end{thinkbox}

\begin{funfact}
世界上最大的油轮之一"诺克·耐维斯号"（已退役），全长458米——比北京从天安门到故宫北门（约960米）的一半还长！它满载时可以装下约56万吨石油。这么大的钢铁巨兽能浮在水上，靠的就是阿基米德原理——它排开的海水量，相当于56个标准游泳池那么多！
\end{funfact}

\newpage

\section{\textcolor{orange!80!black}{第三课\quad 瓦特与烧水壶——蒸汽的力量}}

\subsection*{厨房里的大发现}

两百多年前，在苏格兰的一个小镇上，住着一个叫\textbf{詹姆斯·瓦特}的小男孩。小瓦特特别喜欢观察身边的一切——蚂蚁怎么搬家、钟表怎么走动、风车怎么转动……他对什么都充满了好奇。

瓦特的奶奶家有一个烧水壶，挂在壁炉上烧水。小瓦特经常坐在旁边，看水壶里的水咕嘟咕嘟地沸腾。他发现了一件很有意思的事：水烧开以后，壶盖子会被蒸汽顶得"噗噗噗"地上下跳动，好像有一只无形的手在掀盖子！

小瓦特觉得太神奇了——看不见的水蒸气，怎么有这么大的力气？

\begin{figplaceholder}
此处插入小瓦特坐在壁炉前观察烧水壶的温馨场景插画（见 figure/requirements/nature-intro-ch1.txt 插图5）
\end{figplaceholder}

\subsection*{蒸汽的力气从哪来？}

你可能也见过这个场景。水烧开以后，壶嘴里会冒出白色的水蒸气。这些水蒸气看起来轻飘飘的，但它们的力量可不小！

这个力量从哪里来呢？我们来做一个思想实验：

把一壶冷水放在火上烧。水慢慢变热，里面的小水分子们开始越来越兴奋，动得越来越快——就像下课铃响时同学们冲出教室一样。当温度达到100摄氏度的时候，表面的水分子终于"挣脱"了水面的束缚，变成了看不见的水蒸气，飞到了空气中。

可是，如果水在一个密封的壶里烧呢？蒸汽没地方跑，就会不断地在壶里膨胀、挤压——这就是为什么壶盖子会被顶起来。瓦特看到的"噗噗噗"，其实就是蒸汽在和壶盖子"摔跤"，而且蒸汽赢了！

\subsection*{从烧水壶到蒸汽机}

小瓦特长大后，并没有忘记小时候看到的这个神奇现象。那时候的英国，人们已经开始用蒸汽机来抽矿井里的水了。但当时的蒸汽机又大又笨、效率很低，浪费很多煤。

瓦特看着这些老式蒸汽机，想到了小时候的烧水壶。他花了十几年的时间反复改良，终于发明出了一种效率高出好几倍的\textbf{新型蒸汽机}。

瓦特的蒸汽机，就像是把烧水壶里的蒸汽力量关进了一个"铁盒子"里，然后用这股力量来推动轮子转动。轮子一转，机器就能干活了——抽水、纺织、推动火车和轮船……什么都能做！

瓦特改良的蒸汽机，是人类历史上最重要的发明之一。它直接推动了\textbf{工业革命}——从此以后，人们不再只依靠人力、畜力、风力和水力，而是拥有了强大的蒸汽动力。火车轰隆隆地跑起来了，工厂的机器日夜不停地转起来了，整个世界都跟着加速运转起来。

\begin{thinkbox}
你发现了吗？蒸汽机、火车、工厂……这一切的开始，居然是因为一个小男孩坐在厨房里看烧水壶。所以说，\textbf{好奇心是世界上最好的老师}——一个不起眼的观察，也许就藏着一个改变世界的大秘密。
\end{thinkbox}

\begin{funfact}
为了纪念瓦特的伟大贡献，后来人们用"瓦特"（英文是Watt，简称W）作为\textbf{功率的单位}。你家灯泡上写的"40W"，其实就是"40瓦特"。下次看到灯泡的时候，想想那个看烧水壶的小男孩吧！
\end{funfact}

\begin{figplaceholder}
此处插入瓦特蒸汽机的简易结构示意图（见 figure/requirements/nature-intro-ch1.txt 插图6）
\end{figplaceholder}

\subsection*{两百年后：核电站里的"烧水壶"}

你觉得蒸汽机已经是历史了吗？不！实际上，直到今天，瓦特的"烧水壶"思路仍然在为我们工作——而且是以一种你绝对想不到的方式。

这件事听起来可能有点疯狂：\textbf{全世界核电站发电的核心原理，竟然和瓦特的烧水壶一模一样！}

核电站看起来非常高科技——巨大的混凝土外壳、复杂的控制室、严格的安保。但是，如果你把它的外壳"拆开"，会发现里面的发电流程其实很简单：

\begin{center}
\begin{minipage}{0.9\textwidth}
\centering
\textbf{核燃料发热} $\rightarrow$ \textbf{把水烧开变成蒸汽} $\rightarrow$ \textbf{蒸汽推动汽轮机旋转} $\rightarrow$ \textbf{汽轮机带动发电机发电}
\end{minipage}
\end{center}

看到了吗？中间的"烧水变蒸汽、蒸汽推动轮子转"这一步，和瓦特的蒸汽机是同一个原理！不同的只是"用什么来烧水"——瓦特的时代烧煤，现代核电站用的是\textbf{核裂变}产生的巨大热量。

核裂变是怎么回事呢？简单来说，就是某些特殊元素（比如铀）的原子核被中子撞击后分裂成两半，同时释放出惊人的能量。一公斤铀完全裂变释放的能量，相当于燃烧\textbf{2700吨煤}！用这么小的燃料发出这么多的热，来烧同样一壶"水"，这就是核电站的基本逻辑。

甚至像核潜艇、核动力航母这样的"海上巨兽"，也都在肚子里装着一套"烧水-蒸汽-推动"的系统。它们在深海里悄无声息地航行，能量来源说到底还是一壶"水"。

从两百多年前苏格兰厨房里的一个烧水壶，到今天世界上最先进的能源技术，科学的魅力就在这里——\textbf{大道理常常就藏在小现象里，等着一个有心人去发现它。}

\begin{thinkbox}
中国有很多核电站，比如秦山核电站（浙江）、大亚湾核电站（广东）、田湾核电站（江苏）。它们发的电通过电网送到了千家万户——你正在用的灯、电脑、空调，也许就是靠"烧水"发出来的呢！
\end{thinkbox}

\begin{figplaceholder}
此处插入核电站发电原理简图：核燃料 → 加热水 → 蒸汽 → 汽轮机 → 发电（见 figure/requirements/nature-intro-ch1.txt 插图9）
\end{figplaceholder}

\vspace{0.5cm}

\newpage

\section{\textcolor{orange!80!black}{第四课\quad 小蝌蚪找妈妈——生命的变形记}}

\subsection*{一个经典的故事}

"小蝌蚪，细尾巴，黑身子，大脑袋。出生在池塘，水里来安家……"

你可能听过《小蝌蚪找妈妈》的故事——小蝌蚪们在池塘里游来游去，见到鸭子叫妈妈，见到金鱼叫妈妈，闹了好多笑话，最后终于找到了真正的青蛙妈妈。

这个故事不仅温暖有趣，背后还藏着一个大自然中非常神奇的现象——\textbf{变态发育}。

别被这个名字吓到！"变态发育"不是说发育得很奇怪，而是指：\textbf{一种动物在长大的过程中，身体形态会发生巨大的、根本性的变化。}从一颗卵，到一条"鱼"，再到一只"蛙"——这趟生命旅程，简直是自然界的"变形金刚"。

\subsection*{青蛙的一生：四步大变身}

让我们跟着小蝌蚪，看看它到底经历了什么：

\textbf{第一步：卵}。春天来了，青蛙妈妈在池塘里产下一团团透明的卵，外面包着果冻一样的东西，保护着里面的小生命。如果你蹲在水边仔细看，能看到每个卵里面都有一个黑色的小点——那就是未来的小蝌蚪。

\textbf{第二步：小蝌蚪（用鳃呼吸）}。几天后，卵孵化了，小蝌蚪破卵而出。这时候的它，圆头圆脑，拖着一条细长的尾巴——看起来就像一条小鱼。它用\textbf{鳃}呼吸，就像鱼一样，在水里游来游去，靠啃水草和藻类为生。

\textbf{第三步：长出后腿}。大约过了一个半月，奇妙的事情发生了——蝌蚪的身体两侧鼓出两个小突起，然后慢慢伸长，变成了两条\textbf{后腿}。再过一段时间，前腿也长出来了。与此同时，它的体内也在悄悄发生变化：鳃逐渐退化，\textbf{肺}开始发育。它不再只能用鳃呼吸了，而是要准备到陆地上生活了。

\textbf{第四步：尾巴消失，变成青蛙}。随着四条腿越来越强壮，蝌蚪的尾巴开始一点点变短——不是断掉，而是被身体\textbf{重新吸收}了，营养物质被用来供应腿的生长。最后，尾巴完全消失，鳃也彻底没了，一只小青蛙就这样"变"出来了。它跳上岸，用肺和皮肤一起呼吸，开始了全新的生活。

\begin{figplaceholder}
此处插入青蛙生命周期四阶段示意图：卵 → 蝌蚪 → 长腿蝌蚪 → 青蛙（见 figure/requirements/nature-intro-ch2.txt 插图1）
\end{figplaceholder}

\subsection*{什么是两栖动物？}

青蛙属于\textbf{两栖动物}。两栖动物指的是那些"小时候在水里生活、长大后可以上岸"的动物。它们的名字很有意思：两栖，"栖"就是居住的意思——两个地方都住。

除了青蛙，蟾蜍（癞蛤蟆）、蝾螈、大鲵（娃娃鱼）也是两栖动物。它们都有一个共同的特点：小时候像鱼一样用鳃在水里呼吸，长大后长出肺，可以到陆地上呼吸空气。

但它们的皮肤特别薄，需要保持湿润才能辅助呼吸，所以两栖动物通常不会离水太远。

\begin{funfact}
一只青蛙妈妈一次能产下几百甚至上千颗卵！但其中只有很少一部分能长成青蛙——大部分卵和蝌蚪会被鱼、鸟等天敌吃掉。所以青蛙一次产这么多卵，是大自然确保种群延续的"数量策略"。
\end{funfact}

\begin{thinkbox}
除了青蛙，还有哪些动物小时候和长大后长得完全不一样？提示：蝴蝶（毛毛虫变蝴蝶）、蜻蜓（水虿变蜻蜓）……它们的"变身"过程也属于变态发育哦！你还能想到其他例子吗？
\end{thinkbox}

\newpage

\section{\textcolor{orange!80!black}{第五课\quad 为什么肥皂能洗手？——水的"皮肤"}}

\subsection*{一滴水的秘密}

请你做一个小实验（现在就可以试试）：拿一枚一角的硬币，用手指沾一滴水，轻轻滴在硬币上。慢慢地、一滴一滴地加——你会发现，水滴在硬币上形成了一个鼓鼓的"小水包"，圆滚滚的，像一个小小的水晶球。甚至再多加几滴，水也不会马上流走！

为什么水能"鼓"起来呢？这是因为水的表面有一层看不见的\textbf{"皮肤"}——科学上叫做\textbf{表面张力}。

你看不见它，但它的力量可不小。有些小昆虫（比如\textbf{水黾}，俗称水蜘蛛）就是利用表面张力，在水面上轻盈地滑行，而不沉下去。

\subsection*{表面张力是怎么来的？}

物质都是由极小的粒子——\textbf{分子}组成的。水分子有一个特点：它们喜欢互相手拉手。在水体内部，每个水分子都被周围的水分子从四面八方"拉"着，力量互相抵消了。

但是，在水面最表层的水分子就不一样了——它们上面是空气，没有水分子来"拉"它们。于是，水面上的水分子被下面的水分子拼命往下拽，彼此之间也紧紧抱在一起，就形成了一层\textbf{像橡皮膜一样绷紧的表面}。这就是表面张力。

\begin{figplaceholder}
此处插入表面张力原理示意图：水面水分子互相拉紧，水黾在水面滑行（见 figure/requirements/nature-intro-ch2.txt 插图2）
\end{figplaceholder}

\subsection*{肥皂的"破坏力"}

那么，肥皂是怎么洗手的呢？这和表面张力有什么关系？

其实，肥皂（和洗手液、洗洁精）有一个"超能力"：它们能\textbf{破坏水的表面张力}。

当你用清水洗手的时候，水珠只是在手的表面"滚来滚去"，因为水的表面张力让水珠紧紧地抱成团，不愿意散开去接触脏东西。就像一群手拉手的小朋友围成一圈，外面的人很难挤进去。

但肥皂分子可就厉害了。它长得很特别——\textbf{一头喜欢水（亲水端），另一头讨厌水但喜欢油（疏水端/亲油端）}。当肥皂碰到手上的油污时，亲油端就会插进油污里，把油污团团围住；亲水端则朝外，和水分子拉手。这样一来，油污就被肥皂分子裹挟着，"骗"进了水里，被水一冲就冲走了。

\begin{thinkbox}
洗洁精能洗掉盘子上的油，也是同样的道理。不信下次帮妈妈洗碗的时候观察一下——不加洗洁精，油花漂在水面上怎么都洗不掉；加一滴洗洁精，油花立刻散开了。那就是肥皂分子的"亲油端"在起作用！
\end{thinkbox}

\begin{funfact}
为什么医生和老师总是强调要用肥皂或洗手液洗手，而不是光用水冲？因为水的表面张力让清水很难"抓"住细菌和病毒，但肥皂能破坏表面张力，还能包裹住细菌和油污，把它们彻底冲走。科学实验证明，用肥皂洗手可以将手上的细菌数量减少\textbf{90\%}以上！所以，洗手真的是防病最简单也最有效的方法。
\end{funfact}

\begin{figplaceholder}
此处插入肥皂分子包裹油污/细菌的原理示意图（见 figure/requirements/nature-intro-ch2.txt 插图3）
\end{figplaceholder}

\newpage

\section{\textcolor{orange!80!black}{第六课\quad 指南针为什么指向北？——地球是块大磁铁}}

\subsection*{古代中国的伟大发明}

两千多年前的战国时期，中国人发明了一种神奇的工具——\textbf{司南}。它用天然磁石打磨成勺子的形状，放在光滑的铜盘上。不管你怎么转动它，勺子停下来的时候，勺柄总是坚定地指向南方。

这就是世界上最早的指南针。

后来，人们又发明了更方便携带的指南针：把一根磁化了的铁针穿在灯芯草上，让它浮在水面——铁针会自己转向南北方向。这个精巧的发明，后来传到了阿拉伯和欧洲，帮助了郑和下西洋、哥伦布横渡大西洋……可以说，小小指南针推动了整个人类的\textbf{大航海时代}。

\begin{figplaceholder}
此处插入古代司南和指南针的示意图/照片（见 figure/requirements/nature-intro-ch2.txt 插图4）
\end{figplaceholder}

\subsection*{地球是一块超级大磁铁}

指南针为什么总是指向北方？答案既简单又令人震撼：\textbf{因为地球本身就是一块巨大的磁铁。}

在地球的深处（地核），有大量的铁和镍在高温高压下流动，形成了一个巨大的"天然发电机"，产生了一个贯穿全球的\textbf{磁场}。这个磁场就像在地球内部插了一根看不见的巨型磁铁棒——它的一端在北极附近（地磁南极），另一端在南极附近（地磁北极）。

等一下——"北极附近是地磁南极"？没写错吗？

没错！这就是一个特别容易搞混的地方。磁铁的规律是：\textbf{同名磁极互相排斥，异名磁极互相吸引}（就像你把两块磁铁的N极对在一起，它们会互相推开）。指南针的指针上标着"N"的那一端，其实是指针的\textbf{磁北极}。它指向地球的北方——说明地球北极附近藏着的是\textbf{地磁的南极}，这样才能把指针的N极"吸"过去。

所以是：指南针的N极 → 被吸引 → 指向地球的地理北极附近（那里是地磁的南极）。

绕口令一样？没关系，记住一句话就行：\textbf{地球的肚子里有块大磁铁，它的磁极恰好和地理南北极"反过来"，所以指南针能指方向。}

\begin{figplaceholder}
此处插入地球磁场示意图：地球剖面+磁力线+指南针（见 figure/requirements/nature-intro-ch2.txt 插图5）
\end{figplaceholder}

\subsection*{磁场有什么用？}

地球的磁场不仅帮我们找方向，它还是保护地球生命的\textbf{隐形防护罩}。

太阳不停地向外发射出大量带电粒子——叫做\textbf{太阳风}。如果这些高能粒子直接打到地球表面，会对生物造成很大的伤害。幸运的是，地球磁场把大部分太阳风都"挡"在了外面，把它们导向地球的两极。

当太阳风的带电粒子沿着磁力线进入了极地地区的高空大气时，它们撞击大气中的氧原子和氮原子，就会发出美丽的光芒——这就是\textbf{极光}！北极的叫北极光，南极的叫南极光。极光像是大自然用磁场编织出来的一幅流动的彩色画。

\begin{funfact}
地球上有些动物也能感知磁场！比如候鸟每年迁徙几千公里，方向丝毫不差；海龟出生后会离开海滩游向大洋，几十年后还能游回出生的那片沙滩产卵。科学家认为它们体内有一种"生物指南针"，能感受地球磁场的变化。
\end{funfact}

\begin{thinkbox}
找一块小磁铁和一根缝衣针。用磁铁的一极在针上顺着同一个方向摩擦几十次（不要来回磨），针就被磁化了。然后把它轻轻放在一片树叶上，让树叶浮在水面上——你会发现，针自己转到了南北方向。这就是一个自制的指南针！
\end{thinkbox}

\newpage

\section{\textcolor{orange!80!black}{第七课\quad 富兰克林与风筝——电的"真面目"}}

\subsection*{一场危险的风筝实验}

1752年，美国费城。天空乌云密布，狂风大作，一场雷暴正在逼近。

就在所有人都往屋里躲的时候，一个中年男人却带着他的儿子冲出了家门。他们跑到一片空地上，放起了一只风筝。这个"疯子"，就是\textbf{本杰明·富兰克林}。

富兰克林为什么要在大雷雨天放风筝？他不是在玩游戏——他是在做一个前无古人的危险实验。

富兰克林的风筝上系着一根细细的金属丝，风筝线的末端拴着一把金属钥匙。当雷暴来临、闪电划过天空时，风筝线上的细纤维突然一根根竖了起来——这是电荷聚集的迹象。紧接着，富兰克林把手指靠近了钥匙，只听"噼啪"一声，一个小小的电火花在钥匙和他的手指之间跳了出来！

这一刻，富兰克林证明了：\textbf{天上的闪电，和我们梳头发时产生的静电火花，是同一种东西——都是电！}

\begin{figplaceholder}
此处插入富兰克林雷雨天放风筝实验的场景插画（见 figure/requirements/nature-intro-ch2.txt 插图6）
\end{figplaceholder}

\subsection*{电到底是什么？}

你可能会好奇：电"长"什么样？它是一种液体吗？一种气体吗？

都不是。电的本质，是\textbf{极其微小的带电粒子——电子——的定向流动}。

一切物质都由原子组成，而原子由更小的粒子构成：中心的\textbf{原子核}（带正电）和围绕它旋转的\textbf{电子}（带负电）。正常情况下，电子安安稳稳地绕着原子核转。但在某些条件下——比如摩擦、化学反应、磁场变化——电子可以被"拽"出来，从一个原子跑到另一个原子。

当一大群电子沿着同一个方向跑起来的时候，就形成了\textbf{电流}。电流流过灯丝，灯泡就亮了；电流流过手机里的芯片，就能计算出信息；电流流过电机，风扇就转了。

\subsection*{静电和雷电}

明白了电的本质，我们再回头看富兰克林的风筝。

你肯定遇到过静电——冬天脱毛衣时劈里啪啦响，黑暗中还能看到小火花；用塑料梳子梳头发后，梳子能吸起小纸片。这就是\textbf{静电}：毛衣和头发摩擦，把电子"蹭"走了一部分，导致物体带上了多余的电荷。

雷电的道理也是一样的，只不过规模放大了几亿倍。云层中的冰晶和雨滴互相碰撞、摩擦，使云的不同部分分别带上了大量的正电和负电。当正负电荷积累到一定程度，空气再也挡不住它们了——"轰"！一道闪电劈开天空，巨大的电流在云和地面之间瞬间释放。

一道闪电携带的能量，大约是\textbf{50亿焦耳}——足够让一辆电动汽车跑两万多公里。

\subsection*{避雷针的诞生}

幸运的是，富兰克林没有被雷劈中。（郑重提醒：\textbf{千万千万不能模仿这个实验！}富兰克林活下来纯属运气，后来有模仿者因此丧生。）

但富兰克林从这个实验中获得了一个极其有价值的想法：如果在建筑物的最高处装一根金属杆，再用电线把它和大地连接起来，那么当雷电命中建筑物时，电流就会沿着金属杆和导线安全地流入大地，而不是穿过建筑物造成破坏。

这就是\textbf{避雷针}，富兰克林因为它拯救了无数的生命和财产。

\subsection*{电是好朋友——但要守规矩}

富兰克林的故事告诉我们电是什么，但还有一件非常重要的事，我们必须认真地说一说：\textbf{如何安全地"和电做朋友"。}

电给我们带来了光明、温暖和便利，但电也是一位脾气很大的"朋友"——如果你不尊重它，它就会伤害你。下面这些安全规则，请你一定要记住：

\textbf{第一条：千万不要用湿手碰开关和插座。}水会导电——如果手上有水，电流就可能通过水传到你身上。洗完手、洗完澡，一定要先擦干再碰电器。

\textbf{第二条：不要把金属物品插进插座孔里。}金属是电的"高速公路"——钉子、铁丝、剪刀，通通不能往插座里塞。如果你真的需要插插头，一定要用有绝缘保护的标准插头，而且手要握在塑料部分。

\textbf{第三条：看到破损的电线要远离。}如果看到电线外面的塑料皮破了、里面铜丝露出来了，千万不要碰！要马上告诉大人。破损的电线就像一只"露出牙齿的老虎"。

\textbf{第四条：打雷时不要在户外空旷处逗留。}不要在树下躲雨，不要打伞在空旷地方走（伞骨是金属的），不要放风筝——你不想变成富兰克林！应尽快进入室内或车内。

\textbf{第五条：不要同时在一个插座上插太多电器。}每个插座和电线都有"容量上限"，超过了就会发热，严重时会引起火灾。如果你发现插头或插座摸起来很烫，马上拔掉并告诉大人。

\textbf{第六条：攀爬要远离电力设施。}不要爬电线杆，不要在高压线附近放风筝或玩无人机。高压电不需要直接接触——靠得太近它就能"跳"过来（电弧放电）。

记不住这么多？没关系，记住一句总原则就够了：\textbf{电走它的路，你别给它"搭桥"。}你的身体如果成了电的通路，就会发生危险。

\begin{thinkbox}
下次看到高楼上那根细细的金属杆，你知道它是干什么的了吧？它不是天线，也不是装饰——它是把雷电"请"进大地的安全通道。
\end{thinkbox}

\begin{funfact}
富兰克林不仅是科学家，还是政治家、作家和发明家。他参与了美国《独立宣言》的起草，发明了双焦眼镜和富兰克林炉，还绘制了世界上第一张墨西哥湾暖流图。他给自己写的墓志铭是："印刷工本杰明·富兰克林的遗体……"——这个发明了避雷针、揭开了闪电秘密的人，一生都谦虚地称自己是一个"印刷工"。
\end{funfact}

\begin{figplaceholder}
此处插入避雷针工作原理示意图：闪电→避雷针→导线→大地（见 figure/requirements/nature-intro-ch2.txt 插图7）
\end{figplaceholder}

\vspace{0.5cm}

\newpage

\section{\textcolor{orange!80!black}{第八课\quad 种子是怎么长大的？——植物生长的秘密}}

\subsection*{一颗豆子的"苏醒"}

请你做一个简单的小实验：找一颗干黄豆（或者绿豆、花生），放在一个透明的杯子里，底下垫一张湿纸巾，放在窗台上。每天给它洒一点水，然后——耐心等待。

第一天，豆子好像没什么变化。第二天，它变胖了一点，外皮有了皱纹。第三天……咦？豆子的一端冒出了一个小小的、白白的东西。那是它的\textbf{根}。

再过几天，豆子的另一端顶出了绿色的小芽——那是它的\textbf{茎和叶}。又过了几天，豆瓣（子叶）裂开了，里面的小芽越长越高，展开了嫩绿的叶子，朝着窗外的阳光伸去。

一颗干巴巴的、看起来像"死"了一样的豆子，遇到水、空气和温暖，竟然醒过来了。

\subsection*{种子里面藏了什么？}

如果我们把一颗泡涨了的豆子切开来看，会发现里面藏着下面这些"零件"：

\textbf{种皮}：种子外面那层硬壳，像一件防弹衣，保护里面的"宝宝"不被碰伤、不被虫子咬、不被细菌感染。

\textbf{胚芽}：一个极小的芽尖——将来长成茎和叶。

\textbf{胚根}：又一个极小的芽尖，方向相反——将来长成根。

\textbf{子叶（豆瓣）}：肥肥厚厚的那两大片（豆子最常见的就是它）。它是种子的"便当盒"——里面装满了淀粉和蛋白质，在种子发芽的初期供应营养。等叶子长出来开始光合作用之后，子叶就慢慢枯萎了，因为它已经完成了使命。

所以你看，一颗种子根本不是"死"的。它是在\textbf{休眠}——像一头小熊在冬眠，等着合适的时机醒来。

\begin{figplaceholder}
此处插入种子结构示意图：种皮、胚芽、胚根、子叶剖面图（见 figure/requirements/nature-intro-ch3.txt 插图1）
\end{figplaceholder}

\subsection*{"醒过来"需要什么条件？}

种子发芽需要三个条件：

\textbf{第一，水。}种子在干燥的时候几乎不呼吸——像按了暂停键。一旦吸足了水分，里面的细胞就开始活跃起来，各种酶开始工作，把子叶里储存的营养"转化"成小芽能吸收的糖。

\textbf{第二，合适的温度。}太冷的话，种子继续"睡觉"；太热的话，它会被"煮熟"。每一种植物种子都有自己最喜欢的温度——这就是为什么农民伯伯不是什么时候都能播种。

\textbf{第三，空气（氧气）。}发了芽的种子需要呼吸——就像你跑步要喘气一样。如果把种子完全泡在水里不给空气，它就会"淹死"（沤烂）。

有意思的是——你注意到了吗——种子发芽的前期\textbf{不需要阳光}！因为它还不会光合作用，全靠子叶里的存粮过日子。但一旦叶子展开，阳光就成了必需的"饭票"。

\subsection*{它怎么知道该往哪儿长？}

这件事特别神奇：不管你怎样放豆子——横着放、竖着放、倒着放——根总是往下钻，芽总是往上冒。种子好像"知道"上和下。

是的，植物能感知\textbf{重力}。根会顺着重力方向长（向地性），茎会逆着重力方向长（背地性）。所以根永远去找水和土壤深处，茎永远去找阳光和空气。如果把一棵已经长大了的植物放倒，它的茎会自己弯过来，重新向上长。

还有一种现象你可能也见过：放在窗台上的植物，叶子会朝窗户那一面歪过去。这是因为植物还能感知\textbf{光的方向}（向光性）——哪边光强，就往哪边长。

\begin{thinkbox}
找两颗一样的豆子，一颗种在完全黑暗的盒子里，一颗种在窗台上。过一周比一比——黑暗里的豆芽又细又黄又长（它在拼命"伸"出去找光），窗台上的豆芽又粗又绿又壮。光对植物来说，不只是"能量"，还是"方向信号"。
\end{thinkbox}

\begin{funfact}
世界上最大的种子是\textbf{海椰子}的种子，一颗能重达25公斤——和一个七八岁小孩差不多重！世界上最小的种子是兰花的种子，一粒只有灰尘那么大，一棵兰花蒴果里可以有上百万颗种子。而寿命最长的种子可能是\textbf{莲子}——考古学家在辽宁挖出过1000多年前的古莲子，种下去之后竟然还能发芽开花！
\end{funfact}

\newpage

\section{\textcolor{orange!80!black}{第九课\quad 为什么跑步后会喘气？——呼吸的秘密}}

\subsection*{体育课刚跑完的感觉}

你肯定有过这种体验：体育课上跑了一圈——三百米、四百米——停下来以后，胸口一起一伏，大口大口地喘粗气，"呼——哈——呼——哈——"，心跳得咚咚咚的，脸也红了，额头冒汗。

你有没有想过一个问题：\textbf{为什么跑步的时候身体会要求更多的空气？}

你并没有"决定"要喘气——是身体自己"命令"你喘的。就像口渴了会想喝水、饿了会想吃东西一样，喘气是身体在说："我需要更多氧气！快点！"

\subsection*{细胞里的"小火炉"}

你的身体由几十万亿个\textbf{细胞}组成。每一个细胞就像一个微型的"火炉"——它需要"燃料"（葡萄糖）和"助燃剂"（氧气）来产生能量。

这个"燃烧"过程的科学名称叫\textbf{呼吸作用}。注意——这和肺的"呼吸"（吸气和呼气）是两回事。肺的呼吸是把氧气吸进来、把二氧化碳呼出去；细胞的呼吸作用是氧气进入细胞后，把葡萄糖"烧掉"产生能量，供你走路、跑步、思考、心跳……做一切事情。

简化成一句话就是：

\begin{center}
\textbf{葡萄糖 + 氧气 → 二氧化碳 + 水 + 能量}
\end{center}

当你安静地坐着的时候，身体只需要一点点能量，所以不急不慢地呼吸就够了。但当你开始跑步的时候——肌肉细胞需要大量能量来收缩——"火炉"马上开足马力烧！于是氧气迅速被消耗，二氧化碳（"燃烧"产生的废气）大量堆在血液里。

血液里的二氧化碳一多，大脑的呼吸中枢就被"叫醒"了——它发出紧急通知：\textbf{加快呼吸！加快心跳！}于是你的肺开始疯狂吸氧排碳，心脏开始疯狂泵血，把氧气送到全身，把二氧化碳带走。

这就是为什么跑步后会喘气、心跳加速。

\subsection*{为什么运动员不那么容易喘？}

你可能会注意到，体育老师或者运动员跑步的时候，不太喘，好像很轻松的样子。这是为什么？

因为他们长期训练，\textbf{肺活量}变大了。肺活量就是一次最大吸气后能呼出的最大气量。普通小学生的肺活量大约1500-2500毫升（大概一个大可乐瓶那么多），但游泳运动员可以达到5000-7000毫升——是小学生的两三倍。

肺活量大意味着每次呼吸能吸进更多的氧气，心脏每次跳动能泵出更多的血液，所以不需要像普通人那样又喘又跳地"赶工"。

\begin{thinkbox}
你有没有发现——刚跑完步，马上停下来会更喘；但如果慢慢走一会儿，呼吸反而平复得更快？这是因为慢慢走可以帮助血液继续循环，把堆积的二氧化碳慢慢"运"走。所以体育老师总是说"跑完了别马上坐下，走一走"，是有科学道理的。
\end{thinkbox}

\begin{funfact}
为什么在高原上容易喘？因为高原上空气稀薄，每次吸进去的氧气变少了。身体为了补偿，会制造更多的红细胞来"搬运"氧气。这就是为什么高原训练能提高运动员的成绩——等回到平原，血液里的"氧气搬运工"变多了，体力自然更强了。西藏的牦牛和藏羚羊天生就有大量红细胞，所以能在海拔四五千米的地方健步如飞！
\end{funfact}

\begin{figplaceholder}
此处插入人体呼吸系统简图：肺、气管与气体交换示意（见 figure/requirements/nature-intro-ch3.txt 插图2）
\end{figplaceholder}

\newpage

\section{\textcolor{orange!80!black}{第十课\quad 吃下去的饭去哪了？——消化系统旅行记}}

\subsection*{一口饭的奇幻漂流}

现在，请你想象你正在吃一口米饭。你把它放进嘴里，嚼了嚼，咽了下去……

然后呢？这口饭去哪了？

准备好，我们要跟着这口饭，来一趟身体内部的"奇幻漂流"。

\subsection*{第一站：口腔——"粉碎车间"}

饭一进嘴，你的牙齿就开始工作了——门牙把饭切断，臼齿把它磨碎。与此同时，舌头下面和腮帮子内侧的\textbf{唾液腺}分泌出唾液，把饭粒弄得湿润润的。

唾液可不只是水。它里面有一种叫\textbf{唾液淀粉酶}的东西——专门"攻击"米饭里的淀粉，把它分解成更小的糖分子。这就是为什么你把一口馒头在嘴里多嚼一会儿，会慢慢尝到甜味——淀粉在嘴里就已经开始变成糖了！

所以吃饭要\textbf{细嚼慢咽}，这不只是礼仪——更是在给后面的消化器官"减负"。

\subsection*{第二站：食道——"滑梯通道"}

你把饭咽下去之后，它就进入了一条长约25厘米的肌肉管道——\textbf{食道}。食道像一条有弹性的软管，会做波浪一样的收缩运动，把食物往下推。这个过程叫\textbf{蠕动}。

有意思的是，即使你倒立着吃东西，食物照样能被蠕动推到胃里——重力不是必需的。当然，我们不建议你倒立吃饭。

\subsection*{第三站：胃——"酸液搅拌机"}

食物通过食道，进入了一个像气球一样可以胀大的袋子——\textbf{胃}。胃的内壁会分泌\textbf{胃酸}和消化酶。胃酸的酸性非常强（pH值大约1-2，和汽车电瓶酸液差不多），能把食物中的细菌杀死，还能把蛋白质"拆散"。

胃壁还在不停地收缩搅拌，把食物和胃酸混合成像稀粥一样的东西。这时候那口饭已经面目全非了——变成了一团叫\textbf{食糜}的稀糊糊。

你可能会担心：胃酸这么强，胃自己怎么没被"消化"掉？因为胃内壁有一层厚厚的\textbf{黏液保护层}，像一层"防酸涂料"。但如果这层保护被破坏了（比如吃太多太辣的东西、长期压力大），胃酸就会伤到胃壁——这就是胃疼和胃溃疡的来源。

\subsection*{第四站：小肠——"营养吸收中心"}

食糜从胃出来，进入了本次旅行的主角——\textbf{小肠}。小肠虽然叫"小"肠，但它其实有5-7米长——比你的身高长好几倍！它盘曲在肚子里，像一卷长长的软管。

小肠的内壁布满了密密麻麻的细小突起，叫\textbf{小肠绒毛}。这些绒毛就像毛巾上的小毛毛，大大增加了吸收面积。如果把小肠绒毛全部展平，面积大约有\textbf{一个网球场那么大}！

正是在小肠里，食物中的营养——糖、氨基酸、脂肪酸、维生素——被绒毛吸进了血液，运往全身每一个细胞。

\subsection*{最后一站：大肠——"水回收与打包车间"}

经过小肠吸收后，剩下的东西基本是水分和一些消化不了的纤维。这些残渣进入\textbf{大肠}。大肠的主要工作是把水分回收进身体——所以便便在大肠里越待越干。大肠里还住着上万亿个\textbf{肠道细菌}（别怕，它们是"好细菌"），帮我们分解一些人类自己消化不了的东西，还能合成维生素。

最后，经过大约24到48小时的旅程，那口饭剩下的"废渣"被压缩打包，排出体外。旅程结束。

\begin{thinkbox}
从一口米饭到大便，全程大约走过了8-10米长的通道，耗时一两天。你的身体在这段旅程中像一个超级高效的"食品加工厂"——把米饭拆成零件，营养运送给细胞，废料打包扔掉。你吃下去的每一口饭，身体都在认真对待——所以请你也认真对待你吃的东西。
\end{thinkbox}

\begin{funfact}
胃为什么饿的时候会"咕咕叫"？那不是胃在"说话"，而是胃在空腹状态下进行强有力的收缩运动，挤压里面的气体和少量液体发出的声音。它其实在说："上一批食物已经处理完了，下一批什么时候来？"
\end{funfact}

\begin{figplaceholder}
此处插入人体消化系统旅行示意图：口腔→食道→胃→小肠→大肠（见 figure/requirements/nature-intro-ch3.txt 插图3）
\end{figplaceholder}

\newpage

\section{\textcolor{orange!80!black}{第十一课\quad 秋天树叶为什么变黄变红？}}

\subsection*{北京的秋天是一幅画}

如果你生活在北京，你一定知道——每年十月、十一月，整座城市像被打翻了的颜料盘：钓鱼台的银杏大道一片金黄；香山上的黄栌和枫树红得像火。如果碰上一个大晴天，蓝天白云下金灿灿红彤彤的，美得让人走不动路。

可是你有没有想过：这些树整个夏天不都是绿油油的吗？为什么偏偏到了秋天要"换衣服"？

\subsection*{绿色的秘密——叶绿素}

树叶之所以是绿色，是因为叶子里含有大量的\textbf{叶绿素}。叶绿素是植物进行\textbf{光合作用}的关键分子。它可以吸收太阳光中的红光和蓝光（用来制造养分），偏偏不吸收绿光——绿光被反射了出来，射进我们的眼睛，所以我们看到的叶子就是绿色的。

春夏季日照长、温度高，叶绿素被大量制造出来，老的不停分解，新的不停补充。叶子始终保持着鲜绿的颜色。

\subsection*{秋天来了——叶绿素"下班"了}

到了秋天，白天变短，气温下降。树能"感觉到"这些信号。它开始做一件很重要的事：\textbf{为过冬做准备}。

冬天日照弱，光合作用的效率很低，大片的叶子反而会成为水分的"挥霍者"——叶子每时每刻都在通过气孔蒸发水分（蒸腾作用），而冬天的土壤可能已经冻硬了，根部吸水量大大减少。如果继续保持满树绿叶，树会被"渴死"。

于是树做了一个决定：\textbf{让叶子"退休"！}

它在叶柄和树枝连接的地方长出了一层特殊的细胞（离层），慢慢切断了水和养分的运输通道。叶绿素的生产线停了，旧的叶绿素不断分解却不再补货——绿色就这样逐渐褪去了。

\subsection*{黄色和红色是怎么"露"出来的？}

绿色褪去以后，一直被叶绿素掩盖的其他颜色终于"露脸"了。

\textbf{黄色和橙色}来自于\textbf{类胡萝卜素}。这种色素整个夏天都在叶子里，只不过被叶绿素的绿色盖住了，你看不见它。叶绿素一退场，它们就登台了——这就是银杏和白杨在秋天变黄的原因。

\textbf{红色和紫色}来自于\textbf{花青素}，它和类胡萝卜素不同——它不是一直藏在叶子里的，而是秋天才新制造出来的。当气温降低、阳光充足（尤其是昼暖夜凉的时候），树叶会把残留的糖分转化成花青素，以此来保护叶子在脱落前免受阳光损伤。香山红叶的红，就是这个原理。

所以，一棵树秋天变黄还是变红，取决于它叶子里的"库存"：类胡萝卜素多→变黄（银杏）；花青素制造能力强→变红（枫树、黄栌）；两者都少→直接变褐掉落（杨树、梧桐）。

\begin{thinkbox}
你去过香山吗？为什么山顶的枫叶比山下的更红？因为山顶温差更大、光照更强——这正是刺激花青素制造的最佳条件！
\end{thinkbox}

\begin{funfact}
日本有一种叫"鸡爪槭"的枫树，秋天叶子能红得发紫，几乎像假的一样鲜艳。加拿大国旗上的那片红叶，就是枫叶——加拿大被称为"枫叶之国"，东部秋天的枫林漫山遍野如烈火，是全世界最壮观的秋景之一。
\end{funfact}

\begin{figplaceholder}
此处插入树叶变色的对比示意图：春夏绿叶 vs 秋天黄叶/红叶的色素变化（见 figure/requirements/nature-intro-ch3.txt 插图4）
\end{figplaceholder}

\newpage

\section{\textcolor{orange!80!black}{第十二课\quad 小猫小狗为什么要睡觉？——动物也需要"充电"}}

\subsection*{家里那只"睡神"}

你家里养过猫或者狗吗？如果养过，你一定会注意到：猫一天里大部分时间都在睡觉——有时候趴在沙发上睡，有时候蜷在窗台上睡，有时候四脚朝天地睡。狗也差不多——出门遛弯的时候精神抖擞，一回到家就趴在地上呼呼大睡。

你可能会想：它们怎么这么懒啊？

不是懒。是\textbf{所有动物都需要睡眠}——包括人类。

\subsection*{睡觉的时候，身体在"大扫除"}

睡眠不是"关机"，而是\textbf{切换到了维护模式}。

当你睡着的时候，大脑正在做一件非常重要的事情：\textbf{清理垃圾}。白天大脑在不停运转的过程中会产生代谢废物——尤其是神经系统活动后产生的"蛋白垃圾"。如果这些垃圾堆积多了，大脑就会"卡顿"。睡眠期间，脑细胞之间的空隙会扩大，脑脊液像水流一样冲过这些空间，把垃圾冲走。这就像每天晚上给大脑做一次"大扫除"。

不只大脑在忙——全身都在忙。身体利用睡眠时间修复受损的细胞、巩固免疫系统、分泌生长激素（这就是为什么小孩子需要睡得更多——他们正在长身体）。

\begin{figplaceholder}
此处插入不同动物睡眠时长对比图（条形图/卡通风格）（见 figure/requirements/nature-intro-ch3.txt 插图5）
\end{figplaceholder}

\subsection*{不同动物的"睡眠策略"}

每种动物的睡眠习惯都不一样——这取决于它们在自然界中扮演什么角色。

\textbf{猫：一天睡12-16小时。}猫在野外是顶级捕食者，捕猎需要爆发力，非常消耗能量。所以它们用长时间睡眠来保存能量，为下一次"爆发"做准备。

\textbf{长颈鹿：每天只睡约2小时（还是断断续续的）。}原因很简单——体型太大，如果在草原上躺着睡太久，狮子就会找上门。长颈鹿大多数时候站着打盹，每次只睡几分钟。

\textbf{蝙蝠：一天睡18-20小时，而且是倒挂着睡。}蝙蝠倒挂是因为它们的后腿太弱，不能站在地上；而且一旦遇到危险，爪子一松就能直接起飞。睡这么久是因为它们只在黄昏和夜间活动，飞行消耗能量极大。

\textbf{海豚：大脑只睡一半！}海豚需要定时浮出水面呼吸，如果完全睡着了就会淹死。所以它们进化出了一种神技——\textbf{半脑睡眠}：左脑休息的时候右脑保持清醒，然后换班。你注意过吗？水族馆里的海豚总有一只眼睛是睁着的。

\textbf{人：每天需要8-10小时（小学生要9-11小时）。}如果你睡不够，第二天上课脑子会变"笨"——不是因为你不聪明，而是大脑的"垃圾"还没清理干净。

\begin{thinkbox}
你今晚睡觉前可以做一个"睡眠日记"：记下你几点睡、几点起，第二天上课精神好不好。连续记一个星期，看看睡几个小时你最精神。每个人的最佳睡眠时间都不完全一样——找到你自己的节奏。
\end{thinkbox}

\begin{funfact}
有一种叫\textbf{雨燕}的鸟，可以在飞行中睡觉！科学家发现，雨燕能够在滑翔的时候让大脑的一侧短暂休息，另一侧保持清醒——和海豚一样。雨燕一生中大部分时间都在空中度过，迁徙时能连续飞行10个月不着陆，连睡觉都在天上解决。
\end{funfact}

\vspace{0.5cm}

\newpage

\section{\textcolor{orange!80!black}{第十三课\quad 人为什么会发烧？——身体里的"战争"}}

\subsection*{发烧不是"病"}

你有没有过这样的经历：早上起来头昏昏的，浑身发冷，妈妈一模你的额头——"呀，发烧了！"然后你被裹进被子里，喝着热水，也许还要吃退烧药。

这时候你可能会想：身体怎么这么不争气，偏偏在我明天有重要考试的时候发烧？

其实你冤枉它了。发烧\textbf{不是身体出了问题}，恰恰相反——发烧是身体正在\textbf{努力保护你}的表现。

\subsection*{免疫系统：身体里的"军队"}

要理解发烧，我们先要认识一下身体里的"军队"——\textbf{免疫系统}。

你的身体随时都面临"入侵者"的威胁：病毒（比如流感病毒）、细菌（比如引起喉咙发炎的链球菌），以及其他乱七八糟的病原体。它们无时无刻不想钻进你的身体里繁殖。

但你有一支强大的内部军队在保护你：

\textbf{白细胞（白血球）}是免疫系统的"战士"。它们在血液里巡逻，一旦发现陌生的细菌或病毒，立刻发起攻击——有的直接吞掉敌人（吞噬细胞），有的分泌化学物质"毒杀"敌人，有的记录敌人的长相特征以便下次更快识别（记忆细胞）。

\textbf{抗体}是免疫系统的"通缉令"——一种专门针对特定病原体的蛋白质。病毒一碰到对应的抗体，就像坏人被铐上了手铐，动不了了。

\subsection*{发烧——升高"战场温度"}

那么，发烧和这支"军队"有什么关系呢？

当你身体的免疫系统侦察到大量病原体入侵后，会释放一种叫\textbf{致热因子}的化学信号。这种信号传到大脑里一个叫\textbf{下丘脑}的地方——那是身体的"恒温器"。

正常情况下，恒温器设定在37℃左右。但收到"有敌人入侵！"的信号后，下丘脑主动把设定温度调高了——可能调到38℃、39℃甚至更高。于是你的身体开始"加热"：皮肤血管收缩来减少散热（这就是为什么发烧初期你会觉得冷、发抖——发抖是肌肉快速收缩来产生热量）。

为什么要花能量把自己加热呢？因为：

\textbf{第一，高温能减缓病原体的繁殖。}大多数细菌和病毒在37℃时长势最好，温度一高它们就"蔫"了。

\textbf{第二，高温能加速免疫细胞的行动。}白细胞在稍高温度下移动更快、攻击更猛。

所以发烧其实是免疫系统故意制造的一把"火"，想把敌人烧死。你感觉难受——但敌人更难受。

\subsection*{什么时候需要退烧？}

既然发烧是"好事"，为什么大人有时候还给你吃退烧药呢？

因为温度太高也不行。如果体温持续超过40℃，人体的蛋白质可能开始受损——这就像煮鸡蛋，煮过头就不可逆了。所以一般来说：
\begin{itemize}
  \item 38.5℃以下，多喝水、多休息，让免疫系统自己打仗。
  \item 超过38.5℃，并且人非常难受，可以考虑吃退烧药。退烧药的作用不是"杀敌人"，而是暂时把下丘脑的设定温度调回正常——让你舒服一些，但并没有中断免疫系统的工作。
\end{itemize}

当然，具体要不要吃药、吃什么药，一定要听医生或大人的话——这个千万不能自己决定。

\begin{thinkbox}
你知道吗？冷血动物（比如蜥蜴、蛇）没有内部的"恒温器"。它们生病了怎么办？它们会爬到太阳晒得暖暖的石头上去——用外界的热量来给自己"人工发烧"！这叫\textbf{行为发热}，原理和人发烧一模一样。
\end{thinkbox}

\begin{funfact}
人体正常体温并不是严格的37℃——白天体温略高（下午可达37.2℃），晚上睡眠时体温最低（可能只有36℃出头）。而且不同测量位置温度不一样：口腔≈36.3-37.2℃，腋下≈36-37℃，耳朵≈35.8-38℃。你家用的哪种温度计？下次量体温的时候可以注意一下。
\end{funfact}

\begin{figplaceholder}
此处插入免疫系统工作原理简图：白细胞攻击病原体+下丘脑调高体温（见 figure/requirements/nature-intro-ch4.txt 插图1）
\end{figplaceholder}

\newpage

\section{\textcolor{orange!80!black}{第十四课\quad 面包为什么又软又蓬松？——酵母的"呼吸"}}

\subsection*{馒头里的秘密}

你咬一口刚蒸好的馒头或者面包——软软的、弹弹的，掰开一看，里面全是一小窟窿一小窟窿的，像海绵一样。

如果馒头是实心的，咬起来就像啃橡皮——那得是多痛苦的一件事啊！所以馒头和面包之所以好吃，关键就在于那些\textbf{小窟窿}。

这些窟窿是谁"挖"的？答案是：\textbf{肉眼看不见的小小生物——酵母菌！}

\begin{figplaceholder}
此处插入馒头/面包剖面微距与酵母菌卡通形象图（见 figure/requirements/nature-intro-ch4.txt 插图2）
\end{figplaceholder}

\subsection*{酵母菌是个"贪吃鬼"}

\textbf{酵母菌}是一种单细胞的真菌——是的，它和蘑菇是亲戚，只不过只有一个细胞那么大。

做馒头的人把酵母粉（或老面）混进面团里，再加上一点糖或者面粉本身——酵母菌找到糖以后就像熊找到了蜂蜜，开始大吃特吃。它把糖分子吞进去，在自己的细胞里分解，然后\textbf{吐出两样东西：酒精和二氧化碳气体}。

这个过程叫\textbf{发酵}。发酵其实就是酵母菌在"呼吸"——只是它不需要氧气（这叫无氧呼吸或厌氧呼吸），废气就是二氧化碳。

你可以这样理解：你吃进米饭，呼出二氧化碳；酵母菌"吃"进糖，也呼出二氧化碳。

\subsection*{气体怎么变成窟窿？}

二氧化碳是一种气体，所以当酵母菌在面团里不停地"吐气"时，成千上万的小气泡就在面团里冒了出来。

但面团不是液体——它里面有一种叫\textbf{面筋}的蛋白质网络（面粉加水揉出来的），面筋有弹性，像无数个微型气球，把二氧化碳气体兜住了。气泡越来越多、越来越大，但又跑不掉，于是面团就"发"起来了——变得又大又松软。

最后，当你把发好的面团放进蒸笼或烤箱里加热时，两个事情同时发生：
\begin{enumerate}
  \item 高温把酵母菌杀死了（它们在40℃以上逐渐死亡），发酵停止。
  \item 气泡遇热进一步膨胀，然后面团被蒸/烤"定型"——那些气泡孔就被永久固定下来了。
\end{enumerate}

出锅了。你吃到的每一个小窟窿，都是一个酵母菌留下的"拆迁遗迹"。

\subsection*{苏打和酵母有什么不同？}

你可能还听说过用小苏打发面。小苏打（碳酸氢钠）遇热或遇酸的时候也会分解产生二氧化碳——但它不是生物，只是单纯的化学反应。所以它发面的速度快，但没有酵母发酵过程中产生的特殊香味（酵母除了吐CO₂，还会产生少量酒精和风味物质）。

北方人蒸馒头多用老面（含天然酵母）或酵母粉——这就是为什么手工馒头比泡打粉做的馒头"有面味"。南方有些地方用小苏打+醋来快速发糕，原理一样，但风味不同。

\begin{thinkbox}
你可以和爸爸妈妈一起做一次实验：用三份一样的面团，第一份加酵母，放在温暖的地方（30-35℃），第二份放冰箱里（约4℃），第三份不放酵母。一小时后比比看——第一份胀大了，第二份基本没变（酵母在低温下"懒得动"），第三份完全没变。你亲手验证了发酵需要"酵母+温暖"！
\end{thinkbox}

\begin{funfact}
人类使用酵母发酵的历史至少有五千年以上——古埃及人是最早做面包的文明之一。他们发现把面团放一晚，第二天就"发"起来了，以为是神的恩赐。其实不过是一群看不见的微生物在替他们"吹气球"。今天全世界的面包、馒头、啤酒、酱油、酸奶……全都离不开微生物发酵。你身体里的肠道细菌也在发酵你吃下去的纤维哦——它们帮你，你养它们，互惠互利。
\end{funfact}

\newpage

\section{\textcolor{orange!80!black}{第十五课\quad 可乐里的气泡从哪来？——二氧化碳的"越狱"}}

\subsection*{打开瓶盖的那一刻}

夏天，满头大汗地从外面跑回家，打开冰箱拿出一瓶冰可乐——"啪"一声拧开盖子，"嘶————"，无数小气泡从瓶底涌上来，在液面炸开，喷出凉爽的雾气。仰头喝一大口，气泡在舌头上噼里啪啦地炸裂，那一瞬间，整个世界都清爽了。

这瓶可乐并没有在瓶子里装一台"气泡制造机"。那些气泡其实一直都藏在可乐里——是你把它们放出来的。

\subsection*{高压"囚禁"了二氧化碳}

可乐里的气泡是\textbf{二氧化碳}气体。这种气体在常温常压下是无色的，轻飘飘地飞在空气里。那为什么它能呆在水里？

秘诀有两个词：\textbf{压力和温度}。

在可乐工厂里，工人把二氧化碳在\textbf{高压}下压进水（或糖水）里。在这种高压环境下，二氧化碳分子被"逼"得无处可逃，只好老老实实溶解在水里——就像早晚高峰把你塞进挤满人的公交车上，你虽然不想呆在那里但也出不去。

一旦你拧开瓶盖——"啪"——瓶子里的高压瞬间释放，变成了大气压。那些被"关"了很久的二氧化碳分子终于能"越狱"了！成千上万的小气泡从可乐各处同时冒出来，汇成一股往上冲的白泡泡——这就是你看到的"嘶——"。

\subsection*{温度也有关系}

除了压力，温度也对气泡有很大的影响。

你有没有注意到：冰可乐打开的时候"嘶"一声很克制，温可乐打开的时候泡沫喷涌而出？

这是因为\textbf{气体在冷水里比在热水里更容易溶解}。冰可乐里的二氧化碳"住得比较舒服"，不太急着跑出来；温可乐里的二氧化碳早就想逃跑了，一开门就蜂拥而出——甚至可能喷你一脸。

同样的道理：夏天池塘里的鱼有时会浮到水面张嘴喘气——因为热水里溶解的氧气变少了，水里"闷"得慌，鱼只能浮上来吸大气里的氧气。

\subsection*{为什么摇晃后会喷出来？}

这个问题你大概率亲身经历过。一瓶没开过的可乐掉在地上滚了两圈，然后你打开——"噗——"可乐像喷泉一样喷出来，弄得到处都是。

这是因为摇晃让可乐里的二氧化碳快速聚集成了大量气泡核。正常情况下，溶解在水里的二氧化碳分子是"散兵游勇"——分散在水的各个角落，需要一点时间才能在瓶壁上慢慢聚集成气泡。但摇晃让它们猛烈撞击、瞬间抱团，形成了无数个大气泡。这些大气泡已经蓄势待发，一开瓶——集体冲出来，就喷了。

所以如果你摇了可乐，开瓶前先放一会儿让气泡"冷静"下来，或者用手指弹一弹瓶壁帮助气泡消散——这都是有科学依据的。

\begin{thinkbox}
找一瓶透明汽水（雪碧或七喜最好，因为无色，观察方便），拧开盖子，仔细观察气泡最先从哪里冒出来。你会发现气泡不是均匀冒的——它们大多从瓶壁上或液面附近冒出来。因为气泡的形成需要一个"起点"——瓶壁上的微小瑕疵、液面上的杂质……这就是为什么有时候你往可乐里丢一颗葡萄干，葡萄干表面会疯狂冒泡。
\end{thinkbox}

\begin{funfact}
碳酸饮料是18世纪的英国科学家\textbf{约瑟夫·普里斯特利}偶然发明的。他发现把一碗水放在正在发酵的啤酒桶上方，水里就"充"进了气体，喝起来有刺激的口感。他以为自己发现了"治病神水"，后来被商业开发成了汽水。最早的汽水是在药店里卖的——没错，可乐的前身是"药"！
\end{funfact}

\newpage

\section{\textcolor{orange!80!black}{第十六课\quad 切开的苹果为什么变黑？——氧气的"恶作剧"}}

\subsection*{"变丑"的苹果}

你有没有遇到过这种情况：妈妈切好了一盘苹果，雪白雪白的果肉看着就让人流口水。你出去玩了一会儿回来——咦？苹果变黄了，有的地方甚至变成了褐色，像生锈了一样，看起来完全不像刚才那么好吃了。

你是不是还以为苹果"坏了"？其实没有——它只是和空气发生了一场\textbf{化学反应}。

\subsection*{罪魁祸首：氧气}

这个让苹果变黑的反应叫\textbf{酶促褐变}。让我们把它拆开来看：

\textbf{"酶促"——}苹果果肉里本身就有一种叫\textbf{多酚氧化酶}的化学物质，平时被包在细胞里的一个个小隔间里，老老实实的不乱跑。

\textbf{"褐变"——}当你一刀切下去，刀锋破坏了无数细胞的墙壁，那些小隔间被戳破了。多酚氧化酶被释放出来，遇到了苹果果肉里的另一类物质——\textbf{酚类}（比如绿原酸、儿茶素）。多酚氧化酶立刻"撮合"酚类物质和空气中的\textbf{氧气}发生反应，生成了\textbf{褐色色素（黑色素）}。

所以变黑的本质是：\textbf{切开的细胞 → 酶跑出来 → 酶+酚类+氧气 → 生成褐色物质。}

这和铁生锈是相似的——铁在氧气和水的作用下变成了铁锈（红褐色）；苹果的酚类在酶和氧气的作用下变成了黑色素（黄褐色）。都是"氧化"——只不过苹果的氧化有酶来"帮忙加速"。

\subsection*{怎么救苹果？}

知道了原理，我们就可以想办法"救"苹果了。既然变黑需要——酶、氧气、酚类——三者同时在场，那么阻止其中任何一个环节就行。

\textbf{方法一：滴柠檬汁。}柠檬汁富含维生素C（抗坏血酸），它是一种很强的\textbf{抗氧化剂}——可以抢先和氧气反应，让氧气"没空"去招惹酚类。而且柠檬汁的酸性还能抑制多酚氧化酶的活性。所以滴几滴柠檬汁，苹果就不容易变黑了。这是最常用也最有效的方法。

\textbf{方法二：泡盐水。}淡盐水也能在一定程度上抑制酶的活性，但没有柠檬汁效果好。

\textbf{方法三：隔绝空气。}把切好的苹果用保鲜膜紧紧包起来，或者泡在水里（不让它接触空气），也能延缓变黑。

\begin{figplaceholder}
此处插入苹果氧化对比实验图：未处理 vs 柠檬汁 vs 盐水的对比（见 figure/requirements/nature-intro-ch4.txt 插图3）
\end{figplaceholder}

\subsection*{不只是苹果——氧化无处不在}

苹果变黑只是氧化反应的一个例子。你再看看身边：

\textbf{铁生锈}：铁 + 氧气 + 水 → 铁锈。和苹果变黑同一个机制，只不过没有酶参与。

\textbf{香蕉皮变黑}：也是酶促褐变。所以香蕉不能放冰箱——低温下酶活性更高，反而黑得更快。

\textbf{人的衰老也和氧化有关。}我们身体在利用氧气"燃烧"能量的过程中，会产生一些副产品叫\textbf{自由基}。自由基是一种非常活泼的氧化性物质，会攻击细胞膜和DNA，被认为是衰老和很多疾病的元凶。抗氧化剂（比如维生素C、维生素E、花青素）可以中和自由基，减慢这个过程。所以吃新鲜水果蔬菜，不只是为了维生素——也是在"抗氧化、防衰老"。

\begin{funfact}
有些水果和蔬菜天生就不容易氧化变色——比如柠檬本身含有高浓度维生素C，自带"防腐"功能。而有些植物比如牛油果、土豆，褐变速度极快。食品工业中有专门的"护色剂"（最常见的是维生素C衍生物）加在预先切好的水果里，超市货架上那些切好的苹果放了几天还是白的——不是因为基因好，是化学的功劳。
\end{funfact}

\begin{thinkbox}
下次妈妈切水果的时候，你可以设计一个"小科学实验"：切四片苹果，一片什么都不加（对照组），一片滴柠檬汁，一片泡淡盐水，一片用保鲜膜包起来。等半小时，看看谁最白、谁最黑。你亲手验证了抗氧化和酶促褐变的原理！
\end{thinkbox}

\vspace{0.5cm}

\newpage

\section{\textcolor{orange!80!black}{第十七课\quad 盐是从哪儿来的？——一粒盐的万里路}}

\subsection*{不只是"咸的"}

盐，你可能觉得它就是让菜变咸的调味料。但盐远远不止这个——\textbf{你的身体离不开盐}。

没有盐（准确的说是盐里的钠离子），你的神经就没办法传递信号，你的肌肉就没办法收缩——包括你正在跳动的\textbf{心脏}。你的血液、眼泪、汗液里都有盐。你哭的时候眼泪是咸的，你运动后流的汗也是咸的——身体一直在"提醒"你："我需要盐。"

那问题来了：我们每天吃的盐，到底是从哪儿来的？

盐的来历可以分为三条路：\textbf{海水的路、地下的路、石头的路}。

\subsection*{第一条路：海水晒盐}

这是最古老也最直观的制盐方式。

在海边，人们把海水引进一大片浅浅的池子——叫\textbf{盐田}。太阳晒、海风吹，水慢慢蒸发到空气中（就像湿衣服晾干一样）。水走了，盐留下了——盐是不会蒸发的。随着水分越来越少，水里的盐浓度越来越高，最后盐开始\textbf{结晶}——在水中形成一颗颗微小的立方体盐粒。

如果你用显微镜看盐粒，你会发现它们几乎全是\textbf{小正方体}——这是氯化钠的天然结晶形态，大自然的"几何作业"做得极其工整。

晒出来的盐叫\textbf{海盐}。中国有很长的海岸线，从辽宁到海南，到处都有盐田。天津的长芦盐场是北方最大的海盐产区。

\subsection*{第二条路：钻井取卤}

你可能会想：要是住在离海千里的内陆，盐从哪来？

答案藏在地下。很久很久以前——几千万年甚至上亿年前——四川、云南那些地方本来是海洋。地壳运动让沧海变成了桑田，有一部分海水被"封存"在了地下的岩层里，形成了\textbf{地下卤水}——一种含盐量极高的地下水。

两千多年前的战国时期，四川人就发明了一项世界级技术——\textbf{钻井取卤}。他们在没有现代机械的情况下，用竹子和铁钻头打出了深达百米的井，把地下的卤水抽上来，再用大锅熬干水分，得到盐——这就是\textbf{井盐}。

四川自贡被誉为"千年盐都"。你如果在自贡的燊海井（世界第一口超千米深井，1835年凿成，深1001.42米）参观，会看到古人用竹管从深井里把卤水抽上来的场景——那种技术在当时领先全世界。

\subsection*{第三条路：挖矿取盐}

还有第三种盐——它们根本不在地下的水里，而是\textbf{直接以固体盐层的形式埋在地下}。

这些盐层是怎么来的？亿万年前的海水在一片封闭盆地中完全蒸发，留下了几百米厚的盐层。后来被泥沙覆盖、压成了岩石——这就是\textbf{岩盐}。人们在盐矿里像挖煤一样把岩盐挖出来，粉碎、提纯，就成了我们吃的盐。

欧洲的波兰有一个著名的维利奇卡盐矿，从13世纪开始开采，深入地下一百多米，里面甚至有盐雕的教堂、吊灯和雕像——整个地下宫殿都是用盐建造的。

\begin{figplaceholder}
此处插入三种制盐方式示意图：海水晒盐+钻井取卤+岩盐矿（见 figure/requirements/nature-intro-ch5.txt 插图1）
\end{figplaceholder}

\subsection*{动物也需要盐}

人不是唯一需要盐的生物。山里的山羊会舔岩石上的盐霜；大象会挖含盐的泥土吃；非洲的盐渍地是各种动物的"公共餐厅"——食草动物成群结队去那里舔盐。甚至蝴蝶也会聚集在湿泥地上吸盐分。

所以盐不仅仅是一个调味品——它是\underline{整个生命世界运转不可或缺的化学元素}。

\begin{thinkbox}
你可以做一个小实验验证"蒸发结晶"：在一杯温水里加盐搅拌到不能再溶解为止（饱和盐水）。然后把这杯盐水倒在一个盘子里，放在太阳下晒一两天。水干了以后，盘底会留下一层白白的盐——你亲手晒出了盐！
\end{thinkbox}

\begin{funfact}
人身体里大约含有250克盐——够腌两三条小咸鱼。我们每天通过出汗和排尿流失盐分，所以必须靠食物补充。但如果吃盐太多（每天超过5-6克），血压就会升高，增加心脏病风险。所以盐要"够用但不贪多"——大自然的设计就是如此精妙又需要节制。
\end{funfact}

\newpage

\section{\textcolor{orange!80!black}{第十八课\quad 镜子里的世界——光的反射}}

\subsection*{那个"反着"的你}

每天早上刷牙洗脸的时候，你都会看到一个人——在镜子里。他长得和你一模一样，只是……好像有点不对劲？

举起右手。镜子里的那个人举起了……左手。
眨右眼。镜子里的人眨了左眼。
你往左边走，镜子里的你往右边走。

镜像把你的\textbf{左右}颠倒了。但等一下——你的头顶还在上面，脚还在下面，镜子并没有颠倒\textbf{上下}。这到底是为什么？

\subsection*{光是怎么让你"看到"自己的？}

要回答这个问题，我们要先搞明白镜子是怎么工作的。

你之所以能在镜子里看到自己，是因为\textbf{光的反射}。

光从你身后的灯（或者窗外）发出，照到你的脸上。你的脸把光向四面八方弹开（这叫漫反射）。其中一部分光射到了镜子上——镜子有一层极其光滑的金属镀层（通常是铝或银），会把光整齐地弹回来，而不是散开。

这束从你脸上出发、经镜子弹回你眼睛的光，遵循一个简单的规律——\textbf{反射定律}：

\begin{center}
\textbf{光射到镜子上的角度（入射角）= 光从镜子弹走的角度（反射角）}
\end{center}

就像打乒乓球——球打在球拍上的角度，和弹出去的角度是一样的。

你的大脑"以为"光是沿着直线从镜子里射过来的——于是它在镜子后面"构造"了一个虚像。那个虚像和真实的你大小一样、距离一样，只不过左右颠倒了。为什么左右颠倒了？因为镜子把"前后"方向反了——这一反，左右就跟着反了。

\begin{figplaceholder}
此处插入光的反射定律示意图：入射角=反射角+镜面虚像形成（见 figure/requirements/nature-intro-ch5.txt 插图2）
\end{figplaceholder}

\subsection*{不是所有镜子都是平的}

你见过的镜子大致分三种：

\textbf{平面镜}：你家浴室里的那种。反射的图像和真实物体一样大，只是左右相反。

\textbf{凹面镜（聚光镜）}：镜面像碗一样往里凹。它可以把平行射过来的光汇聚到一个焦点上——温度极高。希腊传说中阿基米德曾用巨大的凹面镜聚焦阳光烧毁了敌军的战船。现代太阳能热发电站用的也是这个原理——用大片凹面镜阵列聚焦太阳光来烧水发电。牙医头上的那个小圆镜也是凹面镜——把灯光聚焦到你嘴巴里。

\textbf{凸面镜（广角镜）}：镜面像球一样往外凸。它把反射的图像缩小，但是\textbf{视野更广}。汽车两侧的后视镜就是凸面镜——所以你看到的后面来车"比实际远一些"。地下车库拐角处的圆形广角镜也是凸面镜，让你看到拐角那边有没有来车。

\textbf{哈哈镜}为什么让人变胖变瘦？因为它的镜面既不平、也不是均匀的弧面——有些地方凸、有些地方凹，不同位置的光反射角度不同，图像就被"拉长"或"压扁"了。

\begin{thinkbox}
站在一面大镜子前，举起右手。镜像里的人举起了哪只手？现在，在镜子上贴一张纸，写上你的名字——镜像里字还正着吗？镜子颠倒的不是左右——是把"前和后"颠倒过来，左右是跟着前后一起被翻的。这是光学里一个经典的"脑筋急转弯"。
\end{thinkbox}

\begin{funfact}
世界上最古老的镜子不是玻璃的——是\textbf{黑曜石}（一种天然火山玻璃）做的。在土耳其发掘出了距今约8000年的黑曜石镜子。中国古代用的是铜镜——把铜磨得可以照人。威尼斯人在15世纪发明了现代意义上的玻璃镜——在玻璃背后镀一层锡和水银的合金，清晰度极高，全欧洲的王室都抢着买。今天的镜子背面镀的是铝，成本低、反射率好。所以每次照镜子的时候，你正在看一项有八千年历史的发明。
\end{funfact}

\newpage

\section{\textcolor{orange!80!black}{第十九课\quad 影子为什么跟着我走？——光的"直线性格"}}

\subsection*{影子的游戏}

你肯定玩过踩影子——在太阳底下跑，后面的朋友追着踩你的影子。你躲来躲去，但影子总是死死地跟着你。

天黑以后，用手电筒照在墙上，用手做出各种造型——小狗、老鹰、小兔子——墙上就出现了对应的"手影"。这就是古老的\textbf{皮影戏}的基本原理。

影子为什么能这么准确地"复制"你身体的轮廓？因为光有一个简单却极其重要的性格：\textbf{光在均匀的介质里沿直线传播}。

\subsection*{光不会"拐弯"}

光是世界上最直的"直性子"。只要前面没有障碍物，光就一直沿直线往前走——不会绕弯、不会跳跃。

如果你的身体挡在了光的前面，光过不去——你身后的地面上就留下了一块黑暗的区域，那就是你的\textbf{影子}。影子的形状，就是你身体"挡住光"的那部分轮廓。

这也能解释为什么影子有时候长、有时候短：

早上，太阳刚升起来，阳光从侧面低低地照过来——你的身体投下的影子拉得老长。中午，太阳在头顶正上方——影子缩在脚下，又短又小。傍晚，太阳又变低了——影子重新拉长。

古希腊人很早就利用这个规律做成了一个"太阳钟"——\textbf{日晷}。日晷上有一根斜立的指针（叫晷针），太阳光照在指针上投下影子，影子落在刻有时辰的盘面上，指到哪个刻度就是什么时间。在钟表发明之前，这是人类计时最精确的工具之一。

\begin{figplaceholder}
此处插入光沿直线传播示意图：早晨长影子→中午短影子→傍晚长影子+日晷原理（见 figure/requirements/nature-intro-ch5.txt 插图3）
\end{figplaceholder}

\subsection*{天体级别的"影子游戏"}

影子不只是地上的小黑块——在宇宙尺度上，影子可以大到遮住整个星球。

\textbf{日食}：月亮飞到了地球和太阳之间，月亮的影子扫过地球表面，站在影子区域里的人看到太阳被"吃掉"了一块（甚至全被遮住）。2020年6月，中国多地观测到了壮观的日环食——太阳变成了一个金色的大戒指。

\textbf{月食}：地球飞到了太阳和月亮之间，地球的影子投射到月亮上。因为地球比月亮大，影子可以完全覆盖月亮——月亮变成暗红铜色（叫"血月"）。这是因为太阳光穿过地球大气层时被折射，红光弯到了月亮上。

日食和月食，就是太阳、地球、月亮三个天体之间庞大的"手影游戏"。

\subsection*{小孔成像——两千年前的"照相机"}

光沿直线传播还有一个神奇的应用：\textbf{小孔成像}。

两千多年前的中国战国时期，\textbf{墨子}和他的弟子们做了一个实验：在一间暗室的墙上开一个小孔，面向外面阳光下的景物。结果发现——对面墙上出现了外面景物的倒立的影像！人如果站在外面，墙上的影子就是头朝下、脚朝上。

这是因为光沿直线传播：景物的顶部发出的光穿过小孔打在墙的下方，底部的光打在墙的上方——上下颠倒了。这就是世界上最早的"照相机"原理，比欧洲早了将近两千年。

你可以自己在家做这个实验：找一个纸杯，底部扎一个小针孔，杯口蒙一层半透明的薄纸（临摹纸）。晚上对着亮着的灯泡——你会发现纸上的灯光影像是倒过来的。

\begin{thinkbox}
用两个手电筒（或者两盏台灯），从不同方向同时照一个物体。你会看到地上出现了\textbf{两个影子}——因为光来自两个不同方向，物体挡住了两条光路。你能想到什么现实场景吗？对了——手术台上的\textbf{无影灯}就是很多盏灯从不同角度同时照过来，把所有影子都互相"覆盖"掉了，医生手术时才不会因为自己的手影挡住视野。
\end{thinkbox}

\begin{funfact}
中国的\textbf{皮影戏}是世界级的非物质文化遗产，有上千年历史。皮影戏用驴皮或牛皮雕刻成各种人物，在幕布后面用灯光照射，观众在幕布前面就能看到彩色的人物"演"故事。皮影戏的原理就是"光沿直线传播→影子"。2011年，中国皮影戏被联合国教科文组织列入人类非物质文化遗产代表作名录。一个小小的光影原理，撑起了一门千年艺术。
\end{funfact}

\newpage

\section{\textcolor{orange!80!black}{第二十课\quad 彩虹是从哪儿来的？——光被"拆开"的魔术}}

\subsection*{"天桥"出现了}

夏天的午后，一场大雨刚停，太阳从云缝里探出脸来。你朝天空另一边看过去——\textbf{一条巨大的彩色拱桥横跨在天上}，从一头到另一头，红、橙、黄、绿、蓝、靛、紫，七种颜色依次排列，像一幅用水彩画在天上的画。

那就是彩虹。古人管它叫"天弓""天桥"，觉得它是神仙走的桥。但科学告诉我们——彩虹不需要神仙，它需要的只是三样东西：\textbf{太阳光、空气中的小水滴、以及你正好站在太阳和水滴之间的某个角度}。

\subsection*{白光不是"纯"的——它是一支"七色队"}

你看到的太阳光（和一般的白光）看起来是"无色"的。但这是假象。

\textbf{白光其实是红、橙、黄、绿、蓝、靛、紫七种颜色光混合在一起的结果。}这就像把七种颜色的橡皮泥揉在一起变成了灰白色——分开来看才知道里面装了什么。

怎么把白光"拆开"呢？需要请一个"拆光高手"出场——\textbf{三棱镜}。

三棱镜是一块三角形的玻璃。当一束白光射进三棱镜的一个面，穿过玻璃从另一个面出来的时候——光被\textbf{折射}了两次（一次进、一次出）。不同颜色的光折射的角度\textbf{不一样}：紫光弯得最厉害，红光弯得最少。于是—束白光从三棱镜里出来后，"哗"地散开了——红橙黄绿蓝靛紫，整整齐齐排成一排。

这就是\textbf{光的色散}。

\subsection*{彩虹就是大自然的三棱镜}

彩虹形成的道理和三棱镜一模一样，只不过三棱镜变成了\textbf{雨后空气中的无数小水滴}。

雨后空气中还悬浮着大量微小的水滴。阳光从你背后射过来（记住——看彩虹的时候太阳一定在你背后），射进远处的水滴：

\begin{enumerate}
  \item 光进入水滴——折射一次（被"掰弯"）
  \item 光在水滴内部碰到内壁——反射回来
  \item 光从水滴里出来——再折射一次
\end{enumerate}

经过"折射→反射→折射"的两次折射，白光被水滴"拆"成了七色光，射进你的眼睛里。无数个水滴都在做同样的事情，但你每个角度看到的光来自不同位置的水滴——紫色来自较低的水滴，红色来自较高的水滴——于是天空中出现了一道弧形的七彩带。

\begin{figplaceholder}
此处插入彩虹形成原理示意图：阳光→水滴→折射+反射+折射→七色光（见 figure/requirements/nature-intro-ch5.txt 插图4）
\end{figplaceholder}

\subsection*{为什么彩虹总是弯的？}

你可能注意到了——彩虹永远是弯的弧线，而不是直的或方的。

这是因为太阳、你的眼睛和你看到的彩虹水滴必须在同一个\textbf{42度的角度关系}上（对于红色光而言，紫色约40度）。所有满足这个角度关系的水滴，在三维空间中构成了一个\textbf{圆锥面}——你站在圆锥的顶点，彩虹就是圆锥的圆形底边。但因为地面挡住了下半截，你只能看到上半截的弧线。

如果你在飞机上，没有地面遮挡——你能看到\textbf{完整的圆形彩虹}！飞行员有时候就能拍到这种奇景。

\subsection*{双彩虹和"霓"}

有时候你能看到彩虹外面还有第二道、颜色更淡、顺序相反的彩虹——外面的彩虹紫色在上、红色在下（和里面那条相反）。这叫\textbf{霓}，也叫\textbf{副虹}。

霓的形成是因为光在水滴里多反射了一次——反射了两次才出来。多一次反射，光就弱了一次，所以霓的颜色比主虹淡。而且因为光路翻转了，颜色的排列顺序也反了过来。

下次看到双彩虹的时候，你可以指着外面那条和朋友说："那不是第二道彩虹，那叫霓——它在水滴里多转了一圈才出来。"

\begin{thinkbox}
你可以在家里"制造"一道彩虹：找一个阳光充足的窗台，放一杯水在窗台上，让阳光透过杯子照在地板上或者白纸上。调整角度——你会在投影的边缘看到一条迷你彩虹！原理一样——水和玻璃替代了水滴和三棱镜的角色。
\end{thinkbox}

\begin{funfact}
CD光盘的背面为什么看上去是彩虹色的？因为光盘表面有极细密的数据轨道（凹槽），间距和光的波长接近。当白光照上去，这些凹槽就起到了类似三棱镜+光栅的作用——把白光分解成了七彩色。所以一片废旧光盘，就是一个可以随身携带的"彩虹制造器"。
\end{funfact}

\vspace{0.5cm}

\newpage

\section{\textcolor{orange!80!black}{第二十一课\quad 温度计为什么能测温度？——热胀冷缩}}

\subsection*{那个"自己会爬"的红色液柱}

你发烧的时候，妈妈把一根细细的玻璃棒夹在你腋下，过几分钟抽出来，眯着眼看上面的刻度："三十八度二，发烧了。"

你有没有想过一个问题：那根玻璃棒里面的红色液体，怎么知道你的身体有多热？

它"知道"，是因为一个几乎所有物质都遵循的自然规律——\textbf{热胀冷缩}。

\subsection*{热胀冷缩——分子们的"舞步"}

所有的物质——固体、液体、气体——都由分子组成。温度是什么？温度其实就是\textbf{分子运动剧烈程度的"评分"}。

温度低的时候，分子们懒懒散散，运动幅度小，彼此靠得近——物质占的体积就小。

温度高的时候，分子们像打了鸡血一样兴奋，运动幅度大，彼此之间的距离被撑开——物质占的体积就大了。

这就是\textbf{热胀冷缩}：温度升高→分子运动加剧→间距增大→体积膨胀。温度降低→分子运动减缓→间距缩小→体积收缩。

温度计里的红色液体（通常是染了红色的酒精或煤油）就是利用这个原理：当它接触到你的热皮肤，液体受热膨胀，体积增大，在细玻璃管中唯一的出路就是往上爬。温度越高，膨胀越多，液柱爬得越高。刻度盘标好了"液柱在这个高度=这个温度"——医生一读就知道你烧了多少度。

\subsection*{水银温度计和酒精温度计}

你可能见过两种温度计：一种是银白色的（水银温度计），一种是红色的（酒精温度计）。它们有什么不同？

\textbf{水银温度计}：里面的液体是汞（水银）。优点是精准、范围广（可测-39℃到356℃）。缺点是水银有毒——一旦打破，汞蒸气会危害健康。现在很多国家已经禁止使用水银温度计了。

\textbf{酒精温度计}：里面的液体是染红了的酒精。优点是安全无毒、能测更低的温度（酒精在-114℃才凝固，所以可以用来测极寒地区的气温）。缺点是不能测太高温度（酒精在78℃就沸腾了）。

你家用的是哪种？现在越来越多的家庭用电子温度计——它不用液体，而是通过热敏电阻来感知温度。

\subsection*{热胀冷缩在生活中无处不在}

温度计是热胀冷缩"最著名"的应用——但绝不是唯一一个。看看你身边：

\textbf{铁轨为什么要留缝隙？}你坐火车的时候，有没有听到"哐当、哐当"的声音？那是火车轮碾过铁轨之间的接缝。铁轨在夏天受热后会变长——如果不留缝隙，铁轨会互相挤压弯曲变形，火车跑上去就危险了。现在的无缝钢轨虽然"无缝"，但也设计了特殊的应力释放结构。

\textbf{为什么夏天的电线变松、冬天的电线变紧？}高压电线在夏天受热膨胀变长，会明显下垂；冬天受冷收缩变短，会绷得紧紧的。架线工人在安装的时候必须计算好，让电线在冬夏之间不会太紧而拉断，也不会太松而碰到树木。

\textbf{为什么玻璃杯倒开水会裂？}玻璃是热的不良导体。往厚玻璃杯里倒开水，杯子的内壁瞬间受热膨胀，但外壁还没热起来、没来得及膨胀——内外膨胀不均匀产生的应力，就把玻璃"撕裂"了。所以倒开水最好用薄杯子，或者先用温水"预热"一下。

\textbf{为什么夏天自行车胎容易爆？}车胎里的空气受热膨胀——如果气打得太满，膨胀到一定程度就会"嘭"。

\begin{figplaceholder}
此处插入热胀冷缩生活场景示意图：温度计+铁轨接缝+电线松紧+玻璃杯开裂（见 figure/requirements/nature-intro-ch6.txt 插图1）
\end{figplaceholder}

\begin{thinkbox}
找一瓶没开过的矿泉水，观察一下瓶子里水的液面在哪个位置。然后把它放进冰箱冷冻层，过一小时拿出来看——液面降低了（水在4℃以上正常热胀冷缩）。但如果冻成了冰——你会发现冰的体积反而比水大了（瓶子胀鼓鼓的，甚至爆裂）。这就是下一课要讲的"水的反常"——水在结冰时不仅不收缩，反而膨胀。大自然总有例外。
\end{thinkbox}

\begin{funfact}
世界上最大的温度计在哪儿？在中国三峡大坝！大坝在建设时，混凝土内部埋设了成千上万个温度传感器——因为巨型混凝土块在凝固时会发热膨胀，如果不监测和控制温度，会产生裂缝。这可能是人类工程史上最大规模的"测温行动"。
\end{funfact}

\newpage

\section{\textcolor{orange!80!black}{第二十二课\quad 为什么先看见闪电后听见打雷？——光与声的"赛跑"}}

\subsection*{雷雨夜的经典疑问}

夏天的夜晚，窗外电闪雷鸣。一道刺眼的闪电劈开黑暗，你下意识地捂住耳朵等——过了几秒，"轰隆隆——"，雷声才滚滚而来。

为什么闪电和雷明明是同时发生的，我们却总是先看到闪电，后听到雷声？

因为\textbf{光和声的速度差得太远了}。

\subsection*{一个快得离谱，一个慢得"像走路"}

\textbf{光的速度}：每秒约30万公里。这个速度有多快？一秒钟，光可以从北京到上海跑300个来回。从地球到月球（约38万公里），光只需要大约1.3秒。你肉眼能看到的任何距离——房间里的灯、窗外的树、天上的飞机——光到达你眼睛的时间基本可以认为是\textbf{瞬间}。

\textbf{声音的速度}：在空气中约每秒340米。这个速度只有光速的\textbf{将近百万分之一}。声音走一公里大约需要3秒。所以当你看到闪电的时候，光几乎零延迟就进了你的眼睛；而雷声要以每秒340米的速度从几公里外"跑"过来——所以你要等好几秒才听得到。

\subsection*{自己动手"量"雷的距离}

这个速度差可以变成一个很好玩的"测距工具"：

\textbf{看到闪电后马上开始数秒——"一千零一、一千零二、一千零三……"——直到听到雷声。每数3秒，雷就大约在1公里之外。}

如果你数到3秒才听到雷声，雷暴大约在1公里外——还算近。如果你数到15秒，那就在5公里外——可以放心地继续站在窗边看。如果你看到闪电后不到1秒就听到了雷声——\textbf{雷暴非常近，赶紧进屋}！

\subsection*{为什么雷声是"轰隆隆"而不是"啪"一声？}

闪电产生的时候，周围空气被瞬间加热到大约3万摄氏度——比太阳表面还热。空气急剧膨胀，产生一道强烈的冲击波——这就是\textbf{雷声}的本质。

但雷声为什么不是清脆的"啪"一声，而是拖长的"轰——隆——隆——"？

有好几个原因：
\begin{enumerate}
  \item 闪电本身很长——一道闪电可能长达几公里，不同部分的雷声到达你耳朵的时间不同，叠成了连续的闷响。
  \item 声音在空气中传播会遇到云层、山体、建筑物的反射——这些回声叠加在一起，让雷声"拖泥带水"。
  \item 声音在远距离传播时，高频部分（尖的声音）衰减得更快，传到远处只剩下低频的"轰隆"声。这就是为什么远处的雷声闷，近处的雷声脆。
\end{enumerate}

\subsection*{光速是宇宙的"限速牌"}

光每秒30万公里的速度，不只是快——它是\textbf{宇宙中最快的速度}。根据爱因斯坦的相对论，任何有质量的东西都不能达到或超过光速。光速是宇宙的"最高限速"。

因为光速是有限的，我们看到的遥远星空其实都是"过去"的影像：太阳光到达地球需要8分20秒，所以你看到的太阳是8分钟以前的太阳。北极星距离我们约430光年——你现在看到的北极星，是明朝时期的北极星发出的光。你仰望星空的时候，其实是在\textbf{看一部宇宙的"历史纪录片"}。

\begin{thinkbox}
下次暴风雨的时候，和爸爸妈妈一起数秒——全家一起"测量"雷暴的距离。看到闪电立刻齐声数"一千零一、一千零二……"，听到雷声立刻停。用秒数除以3，就是雷暴离你的大约公里数。科学、好玩，还锻炼了心算。
\end{thinkbox}

\begin{funfact}
雷声和闪电谁先发生？答案是\textbf{闪电}。雷声是闪电"制造"出来的——闪电加热空气，空气膨胀才产生雷声。所以在同一瞬间是先有闪电后有雷声，只不过两者间隔极短（毫秒级）。再加上光速远快于声速——所以到了你这里，闪电更是"遥遥领先"。
\end{funfact}

\newpage

\section{\textcolor{orange!80!black}{第二十三课\quad 冰为什么能浮在水面上？——水的"反常"性格}}

\subsection*{一块冰的"特立独行"}

在你的印象中，固体是不是应该比液体"沉"？石头丢进水里沉底，铁丢进水里沉底，就连蜡烛丢进水里也沉底。

但冰——固态的水——偏偏反着来：它\textbf{浮}在液态水上。

你喝冰可乐的时候，冰块叮叮当当地浮在最上面；冬天湖面结冰，冰层浮在水面上，冰下面还是液态的水。如果冰会沉底，冬天湖面就不是"结一层冰"了——而是\textbf{从底冻到顶}，整个湖变成一个冰坨坨，湖里所有的鱼虾全得死。

所以冰浮在水面上，不是偶然——是地球生命的幸运。

\subsection*{水在4℃时的"反常"}

大多数物质遵守一个简单规则：\textbf{温度越低，体积越小，密度越大。}

水在大多数温度下也这样——比如100℃的开水比20℃的常温水体积大、密度小。当你把一杯开水放凉，它的体积会略微缩小。

但水有一个"怪脾气"：\textbf{当温度降到4℃以下时，它开始反着来。}

从4℃降到3℃、2℃、1℃……水的体积\textbf{不是继续缩小，而是开始膨胀}。到了0℃结冰的那一瞬间，体积又猛地膨胀了\textbf{大约9\%}——冰的密度变成了约0.9克/立方厘米，而4℃的水密度最大（约1.0克/立方厘米）。因为冰的密度比水小，所以冰浮在了水上。

\subsection*{为什么水会这么"怪"？}

原因出在水分子的排列方式上。

液态水中，水分子可以自由地移动，彼此之间靠得比较近。当温度降低到4℃以下时，水分子之间的\textbf{氢键}开始把分子"锁"进一种更松散、更空旷的结构里——就像一个挤满人的公交车上，大家开始列队站好，反而占了更大的空间。

到了0℃结冰的时候，所有水分子被氢键固定成了一个\textbf{六边形的晶体网格}——就像蜂巢一样，中间有很多空隙。这些空隙让冰的整体体积比同样数量的液态水大了约9\%。

冬天水管被冻爆就是这个原因——水管里的水结冰膨胀，把水管（甚至钢铁管道）给撑裂了，不是冰的力量大，是\textbf{膨胀的空间无处释放}。

\subsection*{"冰浮在水面"救了地球上的生命}

我们现在把目光从一杯冰可乐拉远到整个地球：

冬天，当气温降到零下，湖面开始结冰。冰层浮在水面上，像一床\textbf{被子}，把下面的水和冷空气隔开了。冰层下面的水保持在4℃左右（因为4℃水最重，沉在湖底），不会继续降温——湖底的鱼虾就在这个4℃的"庇护所"里安全过冬。

如果冰沉底，湖面会不断结冰——因为没有冰层在上面遮挡——新结的冰不断沉底，一层一层往上堆，直到整个湖冻实。湖里没有任何生物能活过这样的冬天。同理，如果北冰洋的海冰沉底，全球海洋生态都会崩溃。

你可以说：\textbf{水在结冰时膨胀、冰的密度比水小——这个看似微不足道的"物理反常"，是地球上复杂生命能够存在的基础之一。}

水是地球上最"奇葩"的物质——还有很多反常之处（比如它是少数固态密度小于液态的物质之一）。正是这些"反常"，让地球与众不同。

\begin{thinkbox}
你可以在家做一个观察实验：矿泉水瓶装大半瓶水，标记水位。放进冷冻层，每隔半小时拿出来看一次水位变化。你会先看到水位略微下降（4℃以上的正常热胀冷缩），然后如果继续降温到结冰前，水位会略有回升（4℃以下的反常膨胀）。冻成冰后，冰面明显高于原来的水位——体积增大了！亲手见证水的"怪脾气"。
\end{thinkbox}

\begin{funfact}
为什么溜冰的时候冰刀能滑得那么快？因为冰刀压在冰面上产生了极大的压强，让接触面的一小部分冰瞬间融化成了液态水——形成了一层薄薄的"水膜润滑剂"。冰刀不是"滑在冰上"——而是"滑在水膜上"。这层水膜在冰刀过后会立刻重新结冰，所以你留下的冰痕并不是永久的划伤。物理课本上管这叫"复冰现象"。
\end{funfact}

\newpage

\section{\textcolor{orange!80!black}{第二十四课\quad 恐龙是怎么灭绝的？——"天外来客"改变了一切}}

\subsection*{地球上最著名的"悬案"}

如果让你说出一种已经灭绝的动物，你脱口而出的答案大概率是——\textbf{恐龙}。

恐龙在地球上生活了大约1.6亿年——从2.3亿年前的三叠纪开始，到6600万年前的白垩纪末期突然消失。人类的历史（从早期智人算起）只有大约30万年——连恐龙时代的\textbf{零头}都不到。

这么成功的动物，怎么说没就没了？

科学家为了这个谜题争论了几十年。有人猜是火山大规模喷发、有人猜是气候慢慢变冷、有人猜是哺乳动物偷吃了恐龙蛋……直到一个关键证据被挖了出来。

\subsection*{地层里的"铱"——一封来自太空的信}

1980年，美国物理学家路易斯·阿尔瓦雷兹和他的地质学家儿子沃尔特在意大利的一处地层中发现了一件事：在白垩纪和古近纪的分界线（正好和恐龙灭绝的时期吻合）上，\textbf{铱的含量异常地高}——是正常地层含量的几十倍甚至几百倍。

铱在地球表面非常稀少——它大多沉在地球的核心里。但在小行星和陨石中，铱的含量很高。这意味着什么？意味着在恐龙灭绝的那个时间点，有一颗巨大的\textbf{地外天体}撞击了地球。

阿尔瓦雷兹父子提出了\textbf{小行星撞击假说}——当时被很多科学家嘲笑。但十年后，人们在墨西哥的尤卡坦半岛地下发现了一个直径约180公里的巨型陨石坑——\textbf{希克苏鲁伯陨石坑}——年代正好是6600万年前。证据完全吻合。

\begin{figplaceholder}
此处插入小行星撞击地球导致恐龙灭绝的场景示意图（见 figure/requirements/nature-intro-ch6.txt 插图2）
\end{figplaceholder}

\subsection*{撞击之后发生了什么？}

根据科学家的推算，那颗小行星直径约10公里——和北京五环路围起来的区域差不多大。它以一个极快的速度（每秒约20公里）撞上了地球。

撞击那一瞬间发生了什么？按时间线来看：

\textbf{第0秒：}撞击释放的能量相当于\textbf{几十亿颗广岛原子弹同时爆炸}。撞击点的一切瞬间气化，巨量的岩石、尘土和硫化物被抛向高空。

\textbf{几分钟后：}冲击波引发全球范围的大地震和巨型海啸（海浪可能高达数百米）。

\textbf{几小时到几天后：}被抛上高空的熔岩碎片像流星雨一样落回地面——\textbf{全球"下火雨"}，森林被点燃，大火绵延几大洲。

\textbf{几周到几个月后：}最可怕的来了——撞击激起的尘埃和硫化物气溶胶在高空中扩散，形成了一层厚厚的"天幕"，把太阳光遮住了。地球进入了\textbf{撞击冬天}——全球气温骤降，白天变得像黄昏一样暗。没有阳光，植物无法光合作用，大片死去。吃植物的动物相继饿死，吃肉的动物也因为食物链断裂而灭绝。

\textbf{几年到几十年后：}撞击冬天加上酸雨（硫化物溶解在水里变成硫酸）、臭氧层破坏……大约\textbf{75\%的物种在这次事件中灭绝}。恐龙（非鸟类的全部恐龙）只是其中最著名的一种。

\subsection*{"幸存者"和"接班者"}

但不是所有生物都灭绝了。有些动物活了下来——\textbf{体重小、食性杂、能穴居或能潜水}的生物生存机会更大。

哺乳动物就是最大的"幸存者"之一。恐龙称霸的时代，哺乳动物只能躲着走——大多体型很小（像老鼠那么大），晚上才敢出来活动。恐龙灭绝后，天敌消失了，哺乳动物迅速"占领"了空出来的生态位，开始快速演化、体型增大——最终，其中一支演化成了\textbf{灵长类}，再后来……就是我们现在站在这儿讨论恐龙灭绝的原因了。

如果没有那颗小行星，恐龙可能还在地球上走来走去，而我们人类的祖先可能永远没有机会发展壮大。

\begin{thinkbox}
化石是地球的"历史课本"。在北京，你可以在中国古动物馆看到大量的恐龙化石——包括马门溪龙（脖子超长的中国恐龙）、青岛龙、小盗龙（有羽毛的小型恐龙）等。周末去看看吧——站在巨大的恐龙骨架下面，你会对"6600万年"和"75\%的物种"这些数字有完全不同的感受。
\end{thinkbox}

\begin{funfact}
恐龙并没有"完全灭绝"——有一支活下来了，而且今天还生活在你的身边。那就是\textbf{鸟类}。现代的鸡、麻雀、老鹰、企鹅……它们都是小型兽脚类恐龙的后代。你吃炸鸡的时候，准确地说是在吃"恐龙肉"。科学家在1996年于中国辽宁发现了\textbf{中华龙鸟}化石——一只身上覆盖着原始羽毛的小恐龙，它是"恐龙变鸟"的关键证据。所以下次看到窗外的小麻雀，不妨打个招呼："你好，霸王龙的远房亲戚。"
\end{funfact}

\begin{figplaceholder}
此处插入恐龙灭绝后哺乳动物崛起的时间线简图（见 figure/requirements/nature-intro-ch6.txt 插图3）
\end{figplaceholder}

\vspace{0.5cm}

\newpage

\section{\textcolor{orange!80!black}{第二十五课\quad 天为什么是蓝的？晚霞为什么是红的？——光与大气的"捉迷藏"}}

\subsection*{抬头可见的谜题}

如果你在一个晴朗的秋日中午走出教室，抬头看一眼天空——一大片均匀的、深不见底的蓝色铺满了头顶。但如果到了傍晚，你再朝西方看——蓝色不见了，取而代之的是大片大片的橙红色，"火烧云"把半边天都染透了。

同样的天空、同样的太阳，为什么中午是蓝的、傍晚是红的？

还有一件更奇怪的事：\textbf{太阳光本身是白色的}——这是我们可以确定的（你看正午太阳的"颜色"就是白色刺眼的）。那天空的颜色到底是怎么来的？

\subsection*{瑞利散射——空气分子在"弹"光}

答案来自一位19世纪的英国科学家\textbf{瑞利勋爵}。他发现了一个规律（后来叫\textbf{瑞利散射}）：

\textbf{当光线穿过非常细小的粒子（比如空气分子）时，光会被粒子弹向四面八方。波长短的光（蓝光、紫光）被弹得最厉害；波长长的光（红光、橙光）则比较容易"穿透"过去。}

太阳光进入地球大气层后，遇到无数的氮气分子和氧气分子。这些分子把阳光中的\textbf{蓝光和紫光}拼命往四面八方弹——像一个人拿着霰弹枪把蓝光"散射"到整个天空。

所以，你不管朝天上的哪个方向看，都有被空气分子散射过来的蓝光射进你的眼睛——于是天空看起来是蓝色的。

那紫光去哪了？紫光被散射得比蓝光更厉害，按理说天空应该是紫色的才对。但有两个原因让我们看到的是蓝而非紫：第一，太阳光里紫色本来就比蓝色少；第二，人眼对蓝光比对紫光敏感得多。

\subsection*{为什么傍晚是红的？}

现在我们来看傍晚的天空。

傍晚的时候，太阳在西方地平线附近——你不再"抬头看太阳"，而是太阳在很低很低的角度。阳光不再垂直穿过大气层，而是\textbf{斜着穿过极厚的一层大气}——路程比中午长了十几倍。

在这段漫长的穿越大气的旅途中，蓝光和绿光被散射得"弹尽粮绝"了——还没到达你的眼睛，就已经全部被散射到了其他方向。只剩下\textbf{穿透力最强的红光和橙光}坚持到了最后，直接射进你的眼睛。

于是，你看到的夕阳和晚霞，就是那些"幸存"下来的红光和橙光。云层反射了这些光，就更显壮观——这就是"火烧云"。

\subsection*{顺便说说：星星为什么会"眨眼睛"？}

不知道你有没有想过：为什么星星总是一闪一闪的"眨眼"，而月亮和行星（比如金星、火星）却不闪？

这也是大气层的功劳。星星离我们极远，在天空中只是一个\textbf{点光源}——它的光非常非常细，很容易被大气层中不停流动的气流和温度变化扰动。光线经历了一路的"颠簸"，一会儿偏这边，一会儿偏那边，进入你眼睛的光就时强时弱——看起来就是"一闪一闪"。

而行星（金星、火星、木星）离我们近得多，在天空中是一个极小的\textbf{面}（虽然你肉眼看不出来），有很多条光线同时进入你的眼睛，大气抖动让其中几条光偏掉了，更多条光还在——互相"弥补"，所以整体亮度稳定，不闪。

所以夜空中的一个简单规律是：\textbf{一闪一闪的是恒星，不闪的是行星。}

\begin{figplaceholder}
此处插入瑞利散射示意图：正午蓝天（蓝光散射）+傍晚红霞（红光穿透）（见 figure/requirements/nature-intro-ch7.txt 插图1）
\end{figplaceholder}

\begin{thinkbox}
下次你同时看到金星和旁边的恒星，对比一下——金星完全稳定不闪，旁边的小星星眨啊眨。你能用这个规律把行星从满天星里挑出来！金星（启明星/长庚星）是傍晚或黎明最亮的"星星"，记住——它不闪，因为它是一颗行星。
\end{thinkbox}

\begin{funfact}
如果地球没有大气层，白天你看到的天空不是蓝的——而是\textbf{黑色的}，就像站在月球表面一样。太阳是一个刺眼的白色圆盘挂在漆黑的天空中，星星白天也能看到。阿波罗登月的宇航员就见过这种景象——"正午的天空是黑的"——对当时的人类来说，这是最令人震撼的太空体验之一。
\end{funfact}

\newpage

\section{\textcolor{orange!80!black}{第二十六课\quad 雨是从哪儿来的？——水循环的故事}}

\subsection*{跟着一滴水去旅行}

现在，请你想象自己是一滴海水，躺在一片无边无际的蔚蓝里——比如，在\textbf{渤海}的海面上。

你被太阳晒得暖洋洋的。渐渐地，你身体变得轻盈——不是变轻了，是你和身边的伙伴们开始一个一个地挣脱海面，变成看不见的水蒸气，飞到了空气中。这个过程叫\textbf{蒸发}。

你变成水蒸气以后，和周围的空气混在一起，被风裹挟着往西飞。飞过天津、飞过农田、飞到了\textbf{北京西边的太行山和燕山}脚下。

山挡住了你的去路。你随着气流被迫上升——越往上，温度越低。你从看不见的水蒸气重新凝结成了极小的水滴——和千千万万个和你一样的伙伴聚在一起。这时候，从地面往上看，你们就是一片白白的\textbf{云}。

云里的水滴越来越多、越来越大。气流托不住你们了——你变成了一颗雨滴，\textbf{落了下来}。这就是\textbf{降水}。

你落进了北京密云水库。经过自来水厂的净化，你从水龙头里流出来，被一个四年级的小朋友喝了一口——你又重新进入了生命的循环。或者你落到了泥土地上，渗入地下，流进地下河，最后汇入永定河、海河……奔腾到渤海。你回到了大海，准备下一次旅行。

这就是\textbf{水循环}：蒸发→凝结→降水→径流→再蒸发……一个永不停歇的闭环。

\begin{figplaceholder}
此处插入水循环全景示意图：海洋蒸发→云→降水→河流→海洋（见 figure/requirements/nature-intro-ch7.txt 插图2）
\end{figplaceholder}

\subsection*{为什么迎风坡多雨、北京暴雨是怎么来的？}

你注意到刚才那趟旅程中，山起了什么作用吗？山\textbf{强迫空气上升}。空气一上升就降温，水蒸气就凝结成雨。所以山的迎风坡（风来的那一面）降雨特别多——北京的西南和北部山区降雨就比城区多。

而北京的暴雨，常常是因为\textbf{副热带高压}把大量湿热空气从东南海上推过来，遇到燕山和太行山被迫急剧上升，冷热交汇，形成猛烈的降雨——这就是北京"七下八上"（七月下旬到八月上旬）主汛期的典型天气。

\subsection*{地球是一颗"水球"——但我们缺水}

从太空中看，地球是一颗蓝色的星球——海洋覆盖了\textbf{71\%的地表}。飞船上看起来，地球简直应该叫"水球"。

但如果你以为水是取之不尽的——那就大错特错了。

地球上97\%的水是\textbf{咸水}（海水），不能喝、不能浇地。剩下的3\%淡水里，又有大约70\%被冻在\textbf{南极和格陵兰的冰盖}里，还有一部分在地下的岩层深处——人类真正容易取用的淡水（河流、湖泊、浅层地下水）只占全球总水量的\textbf{不到1\%}。

而中国是一个\textbf{缺水}的国家。华北地区（包括北京）人均水资源量仅为世界平均水平的\textbf{约二十分之一}。北京喝的水有很大一部分来自\textbf{南水北调}——从千里之外的丹江口水库（湖北/河南交界）穿过黄河底部送到北京。

你打开水龙头时哗哗流出的每一滴水，可能是一千多公里之外的一滴长江水穿越了山河，经过了几十道净化工艺，才到了你的杯子里。

\subsection*{节水——不是为了省钱，是为了尊重循环}

理解水循环以后，你就会明白：水不是"用完了就没有了"——它会循环。但人类能从水循环中提取淡水的速度，远远快于水循环补回来。所以节水不是"怕水用完"，而是\textbf{不给水循环系统施加太大的压力}。

具体到每一天：刷牙的时候关掉水龙头（一分钟可以省6升水）、洗澡尽量用淋浴而不是盆浴、洗衣机攒满一桶再洗、看到水龙头滴水马上修好。这些小事，乘以北京两千多万人口，就是一座水库的容量。

\begin{thinkbox}
下一个雨天，拿一个透明杯子放在窗台外面接雨水。等雨停了，看看杯子里的水——它可能来自几百公里外的海洋、经历了蒸发、被风吹过来、从天而降。这就是一滴水的旅行。想一想：这滴水下一次是什么时候才回到大海？也许几天，也许几百年——如果它不幸掉进了冰盖深处，可能要等很久很久。
\end{thinkbox}

\begin{funfact}
你正在喝的水分子，很可能曾经被恐龙喝过——然后又排出来，蒸发、降雨、流入海洋、再蒸发……几十亿年来，地球上的水总量基本没变。你现在身体里60\%的水，每一滴都可能经历过无数次的"前世今生"——曾经是恐龙的眼泪、是秦始皇茶杯里的水、是唐朝的一场暴雨。这真是一个浪漫又科学的想法。
\end{funfact}

\newpage

\section{\textcolor{orange!80!black}{第二十七课\quad 北京为什么春天刮大风沙？——风、沙与看不见的生态链}}

\subsection*{那不是一般的"坏天气"}

如果你在北京生活过几年，你一定经历过这种场景：春天，天空突然变成了土黄色，风裹挟着细沙吹得你睁不开眼，空气中弥漫着一股土腥味。出门走一圈，脸上、头发上、衣服上全是细细的沙土灰。

这就是\textbf{沙尘暴}。北京的孩子对它的感受是"又烦又脏又吓人"——但沙尘暴不仅是一个天气现象，它是一条\textbf{写着因果报应的生态链条}。

\subsection*{沙子从哪儿来？}

北京春天沙尘暴的沙子，主要来自两个地方：\textbf{蒙古高原南部的戈壁和沙漠}，以及\textbf{内蒙古中部的沙地}。

这些地方距离北京多远？最近的浑善达克沙地距离北京直线距离只有180公里——从北京开车往北走三个多小时就到。大风一起，细沙被卷到几千米的高空，顺着西北风一路往东南吹，几个小时后就能覆盖北京。

为什么偏偏是春天？因为春天北方大地刚刚解冻、植被还没完全返青，地表裸露，沙子没有"根"抓。与此同时，冷暖气团交替频繁——冷空气从西伯利亚南下，暖空气从南方北上，交锋产生\textbf{大风}。裸露的地表 + 大风 = 沙尘暴。

\subsection*{沙尘暴背后：一条断掉的生态链}

但沙尘暴最核心的原因不是"风大"——风一直都有。关键是：\textbf{为什么沙子在那里等着被风吹？}

在自然状态下，蒙古高原本来有一层薄薄的植被——草。草根像无数只小手一样紧紧抓着土壤。但是近几十年来，这个地区经历了：

\textbf{过度放牧}：牛羊太多，把草啃得精光。

\textbf{过度开垦}：人们把草原翻成农田，但草原的土壤又薄又贫瘠，种不了几年就废弃了。废弃以后，翻过的土没了草根固着，风一起全飞了。

\textbf{气候变干}：本来就干旱的地区降雨更少了。

于是链条就通了：\textbf{过度放牧和开垦 → 草原退化 → 土地沙漠化 → 地表裸露 → 风一吹 → 沙尘暴 → 吹到北京。}

这就是为什么治理沙尘暴不能只在"风"上下手——你要去\textbf{源头上修复生态}：种树、种草、控制放牧。一棵树、一株草的根，就是一座微型"固沙工程"。

\subsection*{三北防护林——"绿色万里长城"}

为了阻挡风沙，中国从1978年开始了一项人类历史上最大规模的造林工程——\textbf{三北防护林}。它横跨西北、华北、东北（简称"三北"），从新疆一直延伸到黑龙江，规划总面积406.9万平方公里，建设周期超过70年。

这道"绿色长城"的作用不是"把风挡住"——风是挡不住的。它的作用是让地面的植被恢复——\textbf{让沙尘在源头就被"固定在土里"而不是"飞到天上去"。}

经过几十年努力，北京春天的沙尘暴天数已经大幅减少了。但全球气候变化让这个问题变得更加复杂——干旱加剧，沙源地又面临新的压力。所以生态保护不是"一劳永逸"——是需要一代一代人持续投入的事。

\begin{figplaceholder}
此处插入沙尘暴成因与治理示意图：沙源地→风→北京+三北防护林位置（见 figure/requirements/nature-intro-ch7.txt 插图3）
\end{figplaceholder}

\begin{thinkbox}
你在蚂蚁森林里"种"过树吗？支付宝蚂蚁森林上那些虚拟的梭梭树、沙柳，最终会被真实的生态组织种在内蒙古和甘肃的荒漠里。你每一次手机上的点击，对应的是远方土地上多了一棵真正在"抓沙子"的树。一棵梭梭树长大后能固定\textbf{10平方米}的沙地。积少成多，聚木成林——北京的蓝天有一部分就是这些梭梭树"撑"起来的。
\end{thinkbox}

\begin{funfact}
北京在2021年春天遭遇了近十年来最强的一次沙尘暴——PM10浓度一度超过8000微克/立方米（正常值在50以下）。但这次沙尘的沙源地不在内蒙古——它来自距离北京约\textbf{6000公里}之外的\textbf{蒙古国南部}。沙尘暴再次提醒了我们一个深刻的道理：\textbf{生态系统没有国界。}一个国家砍了树，另一个国家就可能吃沙子。生态问题是全人类共同的挑战。
\end{funfact}

\newpage

\section{\textcolor{orange!80!black}{第二十八课\quad 地震是怎么回事？——地球在"活动筋骨"}}

\subsection*{大地为什么会"抖"？}

你如果在北京上学，一定参加过地震演习——警报一响，双手抱头、躲到课桌下面，等"震动"停止后排队疏散到操场。

你有没有想过：\textbf{大地那么坚实，怎么会忽然"抖"起来？}

答案是：大地并没有你以为的那么"安分"。我们脚下的大地，其实一直在\textbf{动}——只是动得极慢，你都感觉不到。

\subsection*{地球的外壳——一副巨大的"活动拼图"}

地球的表面不是一整块硬壳，而是分成了十几块巨大的\textbf{板块}。这些板块不是固定不变的——它们漂浮在地球内部的\textbf{软流层}（像一种极黏稠的热"岩浆糊"）之上，每年移动几厘米。

几厘米听起来不多——和指甲的生长速度差不多。但是\textbf{乘以几百万年}，就会产生翻天覆地的效果：印度板块撞上了欧亚板块，挤出了喜马拉雅山（珠穆朗玛峰现在还在以每年约1厘米的速度长高）；非洲板块正在裂开，几百万年后东非会变成一片新海洋。

板块和板块的交界处，就是地球"脾气最暴躁"的地方。板块之间互相挤压、碰撞、错动——当压力积累到一定程度，岩石突然断裂错位，释放出巨大的能量——\textbf{地震}就发生了。

\begin{figplaceholder}
此处插入地球板块构造示意图：六大板块+地震分布带（见 figure/requirements/nature-intro-ch7.txt 插图4）
\end{figplaceholder}

\subsection*{为什么四川和日本地震多，北京也有风险？}

打开世界地震分布图，你会发现地震的分布不是随机的——它们高度集中在\textbf{板块边界}上。

\textbf{日本}位于四个板块（太平洋板块、菲律宾海板块、欧亚板块、北美板块）交汇的地方——是全世界地震最频繁的国家之一。2011年东日本大地震引发海啸，是很多人记忆犹新的事件。

\textbf{四川}位于印度板块挤压欧亚板块的前沿——龙门山断裂带就在四川盆地和青藏高原交界处。2008年汶川大地震就是这条断裂带上的岩石突然错动。

\textbf{北京}虽然不在板块边界上，但也在地震带上——\textbf{华北地震带}。历史上北京及周边发生过多次强烈地震。1976年\textbf{唐山大地震}（7.8级，距北京约150公里）是20世纪死亡人数最多的地震之一。所以北京的中小学一直坚持地震教育和演习——不是杞人忧天，是\textbf{未雨绸缪}。

\subsection*{地震来了怎么办？——记住三个字}

地震从发生到造成破坏，往往只有十几秒甚至几秒钟。在这短暂的时间里，正确的反应可以救你的命。

记住三个字：\textbf{伏地、遮挡、抓牢}。

\textbf{伏地（蹲下或趴下）}：地震时你可能站不稳——先蹲下来，降低重心，防止摔倒。

\textbf{遮挡（保护头部和颈部）}：躲在坚固的\textbf{桌子下面}——不是桌子旁边，是下面。用一只手护住后颈和头部。如果在教室，就躲在自己的课桌下面。头颈是最脆弱的部分，要优先保护。

\textbf{抓牢（抓住桌腿等固定物）}：一只手护头，另一只手抓住桌腿。地震晃动时桌子可能"跑走"——你要跟着桌子一起动。

还有几个重要的"不要"：

\textbf{不要躲在窗户旁边}——玻璃碎了会飞溅。

\textbf{不要乘坐电梯}——地震时可能断电被困，也怕电梯井道变形卡住。

\textbf{不要站在吊灯、吊扇正下方}——它们可能掉下来。

\textbf{如果在室外}：跑到\textbf{开阔的地方}，远离建筑物、电线杆、广告牌——这些都可能倒塌。

\textbf{破除一个误区}："地震来了站门框下最安全"——这是\textbf{错的}。这种说法源于老式木结构房屋（门框是屋里最坚固的木构件），但现代建筑的承重墙和柱子远比门框结实。站门框下还可能被门扇夹伤。躲桌子下才是正解。

如果你住在高层建筑里，地震时\textbf{不要往外跑}。高楼即使暂时安全，也可能在余震中受损。标准做法是就近躲桌子下，等主震结束后再按疏散路线下楼。

\begin{thinkbox}
和爸爸妈妈一起做一次"家庭地震隐患排查"：看看家里有没有放得不稳的重物（书架顶上、衣柜顶上）？床头挂没挂重画框？电视机放得稳不稳？把这些东西固定好或者移到低处——这就是地震最有效的"预防"。还有一个必做项：全家约定一个震后集合地点。如果地震发生在白天（家人分散在各处），知道去哪儿碰头。
\end{thinkbox}

\begin{funfact}
中国东汉时期的天文学家\textbf{张衡}在公元132年发明了世界上第一台地震仪——\textbf{地动仪}。它像一个青铜酒樽，外面有八条龙，每条龙嘴里含一颗铜珠，对应八个方向。哪个方向发生地震，那个方向的龙嘴就会张开，铜珠落入下面蛤蟆的嘴里。公元138年，西方的龙吐出了珠子——几天后，信使快马赶到，报告陇西（甘肃）发生了地震。这是人类历史上第一次用仪器成功"遥测"远方地震。张衡比西方同类发明早了约1700年。
\end{funfact}

\vspace{0.5cm}

\begin{center}
\rule{0.7\textwidth}{1pt}

\vspace{0.4cm}
{\large\textcolor{blue!70!black}{\textbf{—— 全模块完 ——}}}

\vspace{0.3cm}
\textcolor{gray}{\small 感谢你跟着二十八个小故事一起探索了自然科学的奇妙世界。}\\
\textcolor{gray}{\small 从一颗苹果到宇宙星空，从一滴水到板块运动——希望这些故事在你心里种下了一颗好奇心的种子。}\\
\textcolor{gray}{\small 因为，每一个伟大的科学家，都是从"为什么"开始的。}

\vspace{0.3cm}
\textcolor{gray}{\small 继续问"为什么"吧——这个世界永远等着你去发现。}
\end{center}

\end{document}
